\SourcePage{243}
\SGAsection{5}{伴随群或单连通群的典范分解}
本节将利用第 1 节的结果，把\Term{adjoint-group}{伴随群}（相应地，\Term{simply-connected}{单连通群}）的一个经典\Term{resolution}{分解}推广到\Term{scheme}{概形}情形。为避免行文过于冗长，证明只给出梗概，细节留给读者；实际上，所涉及的始终是\Term{principal-homogeneous-bundle}{主齐性丛}理论中完全标准的证明：结构群的\Term{reduced}{约化}、\Term{twisting}{扭曲}等。

\textbf{5.1。} 回顾（第 XXI 讲 7.4），一个 Dynkin 图是其\Term{connected-component}{连通分支}的\Term{disjoint-union}{不交并}，而这些分支仍为 Dynkin 图。此外，对应于\Term{root-datum}{根数据}的每个\Term{nonempty}{非空}\Term{connected}{连通} Dynkin 图，都\Term{isomorphism}{同构}于第 XXI 讲 7.4.6 列出的某个标准图（$A_n,B_n,\ldots,G_2$）。以下只考虑\Term{connected-component}{连通分支}属于上述类型之一的 Dynkin 图。记这些标准图组成的\Term{set}{集合}为 $\mathsf T$。对任意 Dynkin 图 $D$，记 $n(t)=n_D(t)$ 为 $D$ 中\Term{isomorphism}{同构}于 $t$ 的\Term{connected-component}{连通分支}的个数，其中 $t\in\mathsf T$。$D$ 的类型定义为 $\sum_{t\in\mathsf T}n_D(t)t$。
类型为 $t$ 的 Dynkin 图称为 $t$ 型的\Term{simple}{单图}，类型为 $nt$ 的 Dynkin 图称为 $t$ 型的\Term{isotypic}{等型图}。设 $D_0$ 为 $D$ 的\Term{connected-component}{连通分支}\Term{set}{集合}，设
\[a:D_0\longrightarrow\mathsf T\]
为显然的映射。$D$ 的类型正是 $\sum_{x\in D_0}a(x)$。

\SourcePage{244}
\textbf{5.2。} 设 $S$ 为\Term{scheme}{概形}，$D$ 为满足上述\Term{restriction}{限制}条件的 $S$ 上 Dynkin \Term{scheme}{概形}。\Term{morphism}{态射}对 $L\rightrightarrows D$（$L$ 为 $D$ 的联结\Term{scheme}{概形}）\footnote{编者注：同 3.2 一样，把 $R$ 换成了 $L$。}的\Term{cokernel}{余核}记为 $D_0$。它是 $D$ 的“\Term{connected-component}{连通分支}\Term{scheme}{概形}”（$D$ 为常值时显然存在；一般情形由\Term{descent}{下降}得到）；这是一个\Term{finite}{有限}\Term{twisted-constant}{扭曲常值} $S$-\Term{scheme}{概形}。存在\Term{canonical}{典范}\Term{morphism}{态射}
\[a:D_0\longrightarrow\mathsf T_S.\]
对 $t\in\mathsf T$，置 $a^{-1}(t)=D_{0,t}$；它是 $D_0$ 的子\Term{scheme}{概形}，其在 $D$ 中的\Term{inverse-image}{逆像}记为 $D_t$，称为 Dynkin \Term{scheme}{概形} $D$ 的 $t$ 型\Term{isotypic}{等型分支}。每个 $D_t$ 都是 $D$ 的子\Term{scheme}{概形}，且
\[D=\coprod_{t\in\mathsf T}D_t.\]
注意，若 $D$ 在 $s\in S$ 处的类型为 $\sum_t n(s,t)t$，则 $D_{0,t}$ 在 $s$ 处的次数为 $n(s,t)$。

\textbf{5.3。} 以下只考虑\Term{semisimple}{半单}\Term{adjoint-group}{伴随群}（相应地，\Term{simply-connected}{单连通群}）。为简化表述，我们对\Term{adjoint-group}{伴随群}陈述结果；把所有陈述中的“\Term{adjunction}{伴随}”换成“\Term{simply-connected}{单连通}”，结论仍成立。
回顾，一个\Term{reduced}{约化}\Term{adjunction}{伴随}\Term{root-datum}{根数据}由其 Dynkin 图的类型确定到\Term{isomorphism}{同构}。因而，若\Term{adjunction}{伴随}\Term{root-datum}{根数据} $\mathscr R$ 的 Dynkin 图具有类型 $\sum n(t)t$，就称 $\mathscr R$ 为该类型；若\Term{semisimple}{半单}\Term{adjoint-group}{伴随群} $G$ 的类型由上述类型的\Term{adjunction}{伴随}\Term{root-datum}{根数据}给出，也称 $G$ 具有该类型。若其类型为 $t$（相应地，$nt$，$n>0$），就称 $\mathscr R$ 或 $G$ 为 $t$ 型\Term{simple}{单}（相应地，\Term{isotypic}{等型}）。
若 $G$ 是\Term{semisimple}{半单}\Term{adjoint-group}{伴随群}，则按 5.2 的意义使用记号 $\operatorname{Dyn}_0(G)$ 和 $\operatorname{Dyn}_{0,t}(G)$。

\textbf{5.4。} 设 $t_i$（$i=1,\ldots,n$）为 $\mathsf T$ 中互不相同的元素，$\mathscr R_i$ 为带\Term{pinning}{钉扎}的 $t_i$ 型\Term{isotypic}{等型}\Term{adjunction}{伴随}\Term{root-datum}{根数据}。考虑乘积 $\mathscr R=\mathscr R_1\times\cdots\times\mathscr R_n$（第 XXI 讲 6.4.1）。设 $S$ 为使各个 $G_S^{\mathrm{Ép}}(\mathscr R_i)$ 都存在的\Term{scheme}{概形}（参见第 XXV 讲）。立即可验证存在\Term{canonical-isomorphism}{典范同构}
\[G_S^{\mathrm{Ép}}(\mathscr R)=G_S^{\mathrm{Ép}}(\mathscr R_1)\mathop{\times}_S\cdots\mathop{\times}_S G_S^{\mathrm{Ép}}(\mathscr R_n).\tag{*}\]
此外，若 $E_i$ 表示 $\mathscr R_i$ 的\Term{automorphism-group}{自同构群}，$E$ 表示 $\mathscr R$ 的\Term{automorphism-group}{自同构群}，$D_i$（相应地，$D$）为 $\mathscr R_i$（相应地，$\mathscr R$）的 Dynkin 图，则有\Term{isomorphism}{同构}
\[E_i\simeq\operatorname{Aut}(D_i),\qquad D=\coprod_iD_i,\qquad E\simeq\prod_iE_i\simeq\operatorname{Aut}(D).\]
结合 (*) 与 1.4，可见 (*) \Term{induction-isomorphism}{诱导同构}
\[A_S(\mathscr R)\simeq A_S(\mathscr R_1)\mathop{\times}_S\cdots\mathop{\times}_S A_S(\mathscr R_n).\]

\SourcePage{245}
\begin{SGAresult}{命题}{5.5}
设 $S$ 为\Term{scheme}{概形}，$G$ 为\Term{semisimple}{半单}\Term{adjoint-group}{伴随} $S$-群。存在唯一\Term{resolution}{分解}
\[G\simeq\prod_{t\in\mathsf T}G_t,\]
其中 $G_t$ 是 $t$ 型\Term{isotypic}{等型}\Term{semisimple}{半单}\Term{adjoint-group}{伴随} $S$-群。此外，该\Term{resolution}{分解}\Term{induction-isomorphism}{诱导同构}
\[\operatorname{Dyn}_t(G)\simeq\operatorname{Dyn}(G_t),\qquad \operatorname{Aut}_{S\text{-gr.}}(G)\simeq\prod_{t\in\mathsf T}\operatorname{Aut}_{S\text{-gr.}}(G_t).\]
\end{SGAresult}
当 $G$ \Term{split-group}{分裂}时，上面已证明此结论。一般情形可假设 $G$ 具有常值类型 $\mathscr R$。利用 $A_S(\mathscr R)$ 的上述\Term{resolution}{分解}及 1.17，就得到所需的 $G$ 的\Term{resolution}{分解}。其余结论由\Term{descent}{下降}证明。

\textbf{注记 5.6。} 更一般地，若 $G,H$ 为两个\Term{semisimple}{半单}\Term{adjoint-group}{伴随群}，则有如下\Term{canonical-isomorphism}{典范同构}（图表交换）：
\[
\begin{tikzcd}[column sep=large]
\operatorname{Isom}_{S\text{-gr.}}(G,H) \arrow[r,"\sim"]\arrow[d]&\prod_{t\in\mathsf T}\operatorname{Isom}_{S\text{-gr.}}(G_t,H_t)\arrow[d]\\
\operatorname{Isomext}(G,H)\arrow[r,"\sim"]\arrow[d,"\sim"']&\prod_{t\in\mathsf T}\operatorname{Isomext}(G_t,H_t)\arrow[d,"\sim"]\\
\operatorname{Isom}_{\mathrm{Dyn}}(\operatorname{Dyn}(G),\operatorname{Dyn}(H))\arrow[r,"\sim"]&\prod_{t\in\mathsf T}\operatorname{Isom}_{\mathrm{Dyn}}(\operatorname{Dyn}(G_t),\operatorname{Dyn}(H_t)).
\end{tikzcd}
\]
\textbf{注记 5.7。} 可以给出 $G_t$ 的内在刻画，这里只陈述而不证明：它是 $G$ 中最大的\Term{normal}{正规}（而且是\Term{characteristic}{特征}的）、$t$ 型\Term{isotypic}{等型}\Term{reductive-group}{约化子群}。
上述命题把\Term{semisimple}{半单}\Term{adjoint-group}{伴随群}的研究归结为\Term{isotypic}{等型}\Term{semisimple}{半单}\Term{adjoint-group}{伴随群}的研究。以下就进行这项研究。

\textbf{5.8。} 设 $\mathscr R$ 为 $t$ 型\Term{simple}{单}、\Term{reduced}{约化}、带\Term{pinning}{钉扎}的\Term{adjunction}{伴随}\Term{root-datum}{根数据}，$I$ 为\Term{nonempty}{非空}\Term{finite}{有限}\Term{set}{集合}，$\mathscr R^I$ 为 $\mathscr R$ 的各个副本 $\mathscr R_i$（$i\in I$）的乘积\Term{root-datum}{根数据}。记 $E$ 为 $\mathscr R$ 的\Term{automorphism-group}{自同构群}；它等同于 $\mathscr R$ 的 Dynkin 图 $D$ 的\Term{automorphism-group}{自同构群}。$\mathscr R^I$ 的 Dynkin 图等同于 $D^I$，因而有\Term{exact-sequence}{正合列}
\[1\longrightarrow\operatorname{Aut}(D)^I\longrightarrow\operatorname{Aut}(D^I)\longrightarrow\operatorname{Aut}(I)\longrightarrow1,\]
其中 $\operatorname{Aut}(I)$ 表示 $I$ 的\Term{permutation-group}{置换群}。因此，\Term{canonical-isomorphism}{典范同构}
\[G_S^{\mathrm{Ép}}(\mathscr R)^I\simeq G_S^{\mathrm{Ép}}(\mathscr R^I)\]
\SourcePage{246}
诱导\Term{exact-sequence}{正合列}
\[1\longrightarrow A_S(\mathscr R)^I\longrightarrow A_S(\mathscr R^I)\longrightarrow\operatorname{Aut}(I_S)\longrightarrow1,\]
最后一项是对应于 $\operatorname{Aut}(I)$ 的常值 $S$-群。另一方面，$I$ \Term{canonical}{典范}地等同于 $D^I$ 的\Term{connected-component}{连通分支}\Term{set}{集合}。若 $G$ 为 $\mathscr R^I$ 型\Term{semisimple}{半单} $S$-群，并定义了一个 $A_S(\mathscr R^I)$ 下的\Term{principal-homogeneous-bundle}{主齐性丛} $P$（参见 1.17），则 $\operatorname{Dyn}(G)$ 的\Term{descent}{下降}定义（3.7）与 $\operatorname{Dyn}_0(G)$ 的定义（5.2）表明，$P$ 通过\Term{morphism}{态射} $A_S(\mathscr R^I)\to\operatorname{Aut}(I_S)$ 关联的丛，正是对应于 $I_S$ 的\Term{form}{形式} $\operatorname{Dyn}_0(G)$ 的 $\operatorname{Isom}_S(I_S,\operatorname{Dyn}_0(G))$。再次利用\Term{category-equivalence}{范畴等价} 1.17 与上述\Term{exact-sequence}{正合列}，通过\Term{form}{形式}推理可得，存在一个 $\mathscr R$ 型\Term{reductive-group}{约化} $\operatorname{Dyn}_0(G)$-群 $G_0$，以及一个 $S$-\Term{isomorphism}{同构} $G\simeq\prod_{\operatorname{Dyn}_0(G)/S}G_0$。

\begin{SGAresult}{命题}{5.9}
设 $S$ 为\Term{scheme}{概形}，$G$ 为 $t$ 型\Term{isotypic}{等型}\Term{semisimple}{半单}\Term{adjoint-group}{伴随} $S$-群。存在（唯一的）$t$ 型\Term{simple}{单}\Term{semisimple}{半单}\Term{adjoint-group}{伴随} $\operatorname{Dyn}_0(G)$-群 $G_0$ 和 $S$-\Term{isomorphism}{同构}
\[G\simeq\prod_{\operatorname{Dyn}_0(G)/S}G_0.\]
此外，此\Term{isomorphism}{同构}诱导\Term{exact-sequence}{正合列}
\[1\longrightarrow\prod_{\operatorname{Dyn}_0(G)/S}\operatorname{Aut}_{S\text{-gr.}}(G_0)\longrightarrow\operatorname{Aut}_{S\text{-gr.}}(G)\longrightarrow\operatorname{Aut}_S(\operatorname{Dyn}_0(G))\longrightarrow1.\]
\end{SGAresult}
显然可以假设 $G$ 具有常值类型 $nt$。第一条已获证明（唯一性显然）。第二条由\Term{split-group}{分裂}情形通过\Term{descent}{下降}得到。
可以合并 5.5 与 5.9：
\begin{SGAresult}{命题}{5.10}
设 $S$ 为\Term{scheme}{概形}，$G$ 为\Term{semisimple}{半单}\Term{adjoint-group}{伴随} $S$-群，$\mathscr D=\operatorname{Dyn}(G)$ 为其 Dynkin \Term{scheme}{概形}。

(i) 存在\Term{canonical}{典范}\Term{resolution}{分解}
\[G\simeq\prod_{t\in\mathsf T}\prod_{\mathscr D_{0,t}/S}G_{0,t}\simeq\prod_{\operatorname{Dyn}_0(G)/S}G_0,\]
其中每个 $G_{0,t}$ 都是 $t$ 型\Term{simple}{单}\Term{adjoint-group}{伴随} $\operatorname{Dyn}_{0,t}(G)$-群（相应地，$G_0$ 为\Term{semisimple}{半单}\Term{adjoint-group}{伴随} $\operatorname{Dyn}_0(G)$-群，其在 $x\in\operatorname{Dyn}_0(G)$ 处的类型为 $a(x)\in\mathsf T$；\Term{morphism}{态射} $a:\operatorname{Dyn}_0(G)\to\mathsf T_S$ 已在 5.2 定义）。

(ii) 上述\Term{resolution}{分解}诱导如下\Term{isomorphism}{同构}的\Term{exact-sequence}{正合列}（竖列），其中 $\operatorname{Aut}_{S,a}(\mathscr D_0)$ 表示与\Term{morphism}{态射} $a:\operatorname{Dyn}_0(G)\to\mathsf T_S$ 交换的 $\operatorname{Dyn}_0(G)$ 的\Term{automorphism}{自同构}\Term{scheme}{概形}：
\SourcePage{247}
\[
\begin{tikzcd}[column sep=large]
1\arrow[d]&1\arrow[d]\\
\prod_{t\in\mathsf T}\prod_{\mathscr D_{0,t}/S}\operatorname{Aut}_{\mathrm{gr.}}(G_{0,t})\arrow[d]\arrow[r,"\sim"]&\prod_{\operatorname{Dyn}_0(G)/S}\operatorname{Aut}_{\mathrm{gr.}}(G_0)\arrow[d]\\
\operatorname{Aut}_{S\text{-gr.}}(G)\arrow[d]\arrow[r,equal]&\operatorname{Aut}_{S\text{-gr.}}(G)\arrow[d]\\
\prod_{t\in\mathsf T}\operatorname{Aut}_S(\operatorname{Dyn}_{0,t}(G))\arrow[r,"\sim"]\arrow[d]&\operatorname{Aut}_{S,a}(\operatorname{Dyn}_0(G))\arrow[d]\\
1&1
\end{tikzcd}
\]
\end{SGAresult}
\begin{SGAresult}{推论}{5.11}
在上述条件下，以下三个\Term{category-equivalence}{范畴等价}：

(i) $G$ 下的\Term{principal-homogeneous-bundle}{主齐性丛}\Term{category}{范畴}；

(ii) $G_0$ 下的\Term{principal-homogeneous-bundle}{主齐性丛}\Term{category}{范畴}；

(iii) 对 $t\in\mathsf T$，各 $G_{0,t}$ 下的\Term{principal-homogeneous-bundle}{主齐性丛}\Term{category}{范畴}的\Term{product-category}{乘积范畴}。
\end{SGAresult}
这由上述\Term{resolution}{分解}和 8.4 \Term{form}{形式}地得到。
\begin{SGAresult}{推论}{5.12}
有\Term{canonical-isomorphism}{典范同构}
\[\operatorname{Tor}(G)\simeq\prod_{t\in\mathsf T}\prod_{\mathscr D_{0,t}/S}\operatorname{Tor}(G_{0,t})\simeq\prod_{\operatorname{Dyn}_0(G)/S}\operatorname{Tor}(G_0),\]
把 $\operatorname{Tor}$ 换成 $\operatorname{Bor}$（相应地，$\operatorname{Kil}$）也一样。
\end{SGAresult}
由\Term{split-group}{分裂}情形立即可得。
\textbf{注记 5.13。} \Term{canonical}{典范}\Term{morphism}{态射}
\[\operatorname{Dyn}(G)\longrightarrow\operatorname{Dyn}_0(G)\]
使我们可以把 $\operatorname{Dyn}(G)$ 看作 $\operatorname{Dyn}_0(G)$ 上的 Dynkin \Term{scheme}{概形}；它实际上就是 $G_0$ 的 Dynkin \Term{scheme}{概形} $\operatorname{Dyn}(G_0)$。
同样，若 $T\subset B$ 为 $G$ 的 Killing 对，\Term{canonical}{典范}地对应于 $G_0$ 的 Killing 对 $T_0\subset B_0$，则可验证，$G$ 与 $G_0$ 的\Term{quasi-split}{拟分裂}障碍相同；这些障碍位于（3.9）$\operatorname{Pic}(\operatorname{Dyn}(G))=\operatorname{Pic}(\operatorname{Dyn}(G_0))$ 中。因此得到：
\begin{SGAresult}{推论}{5.14}
以下条件\Term{equivalence}{等价}：

(i) $G$ \Term{quasi-split}{可拟分裂}；
\SourcePage{248}
(ii) $G_0$ \Term{quasi-split}{可拟分裂}；

(iii) 每个 $G_{0,t}$（$t\in\mathsf T$）都\Term{quasi-split}{可拟分裂}。
\end{SGAresult}

\SGAsection{6}{约化群的 Borel 子群的自同构}
\begin{SGAresult}{引理}{6.1}
设 $S$ 为\Term{scheme}{概形}，$(G,T,M,R)$ 为\Term{split-group}{分裂} $S$-群，$\Delta$ 为 $R$ 的\Term{simple-root-system}{单根系}，$B$ 为对应的 \Term{borel-subgroup}{Borel 子群}。则 $B^u$ 作为 (fppf) 群\Term{sheaf}{层}由 $U_\alpha$（$\alpha\in\Delta$）生成。
\end{SGAresult}
设 $H$ 为由 $U_\alpha$（$\alpha\in\Delta$）生成的 $B^u$ 的\Term{subgroup}{子群}\Term{sheaf}{层}。对整数 $\operatorname{ord}(\beta)=\operatorname{ord}_\Delta(\beta)>0$（参见第 XXI 讲 3.2.15）归纳，证明 $H\supset U_\beta$（$\beta\in R_+$）。当 $\operatorname{ord}(\beta)=1$ 时，由假设成立。因而设 $\operatorname{ord}(\beta)>1$，且只要 $\operatorname{ord}(\gamma)<\operatorname{ord}(\beta)$ 就有 $U_\gamma\subset H$。存在 $\alpha\in\Delta$ 使 $\beta-\alpha\in R_+$（第 XXI 讲 3.1.1）。设 $p$ 为使 $\beta-p\alpha=\beta'\in R_+$ 的最大整数。有 $U_\alpha,U_{\beta'}\subset H$，且 $\beta'-\alpha\notin R$。因此归结为证明：
\begin{SGAresult}{引理}{6.2}
沿用第 XXIII 讲 6.4 的记号。设 $p=1$，即 $\beta-\alpha$ 不是根。若 $H$ 为 $B^u$ 的\Term{subgroup}{子群}\Term{sheaf}{层}，且 $U_\alpha,U_\beta\subset H$，则只要 $i\alpha+j\beta\in R$，$i>0,j>0$，就有 $U_{i\alpha+j\beta}\subset H$。
\end{SGAresult}
按 $q=0,1,2,3$ 的取值分为四种情形。以下 $x,y$ 表示 $\mathbb G_{a,S'}$ 的任意两个\Term{section}{截面}，其中 $S'\to S$。
若 $q=0$，无须证明。
若 $q=1$，则
\[p_{\alpha+\beta}(\pm x)=p_\beta(-y)p_\alpha(-x)p_\beta(y)p_\alpha(x)\in H(S'),\]
故 $U_{\alpha+\beta}\subset H$。
若 $q=2$，同样有
\[p_{\alpha+\beta}(\pm xy)p_{2\alpha+\beta}(\pm x^2y)=p_\beta(-y)p_\alpha(-x)p_\beta(y)p_\alpha(x)\in H(S'),\]
因此适当改变某些\Term{signature}{符号}可得
\[F(x,y)=p_{\alpha+\beta}(xy)p_{2\alpha+\beta}(x^2y)\in H(S').\]
因为 $p_{\alpha+\beta}$ 与 $p_{2\alpha+\beta}$ 交换，故
\[F(x,1)F(1,-x)=p_{2\alpha+\beta}(x^2-x).\]
\footnote{编者注：把 $F(-1,x)$ 改正为 $F(1,-x)$。}
若 $a\in\mathbb G_a(S)$，方程 $X^2-X=a$ 定义 $S$ 的一个自由且\Term{surjectivity}{满射}的\Term{group-extension}{扩张}（即 $\operatorname{Spec}\mathscr O_S[X]/(X^2-X-a)$）；所以 $p_{2\alpha+\beta}(a)\in H(S')$，故 $U_{2\alpha+\beta}\subset H$，进而也有 $U_{\alpha+\beta}\subset H$（因为 $F(1,a)\in H(S')$）。
若 $q=3$，同样有
\[F(x,y)=p_{\alpha+\beta}(xy)p_{2\alpha+\beta}(x^2y)p_{3\alpha+\beta}(x^3y)p_{3\alpha+2\beta}(x^3y^2)\in H(S');\]
又因为
\[p_{3\alpha+2\beta}(\pm x)=F(1,1)^{-1}p_\beta(-x)F(1,1)p_\beta(x)\in H(S'),\]
得到 $U_{3\alpha+2\beta}\subset H$，因而
\[K(x,y)=p_{\alpha+\beta}(xy)p_{2\alpha+\beta}(x^2y)p_{3\alpha+\beta}(x^3y)\in H(S').\]
现在模 $U_{3\alpha+2\beta}$ 计算
\[K(x+y,1)K(-x,1)K(1,y)^{-1},\]
\SourcePage{249}
得到
\[p_{2\alpha+\beta}(2x^2+y^2+2xy-y)p_{3\alpha+\beta}(y^3+3xy^2+3x^2y-y)\in H(S').\]
若 $a\in\mathbb G_a(S)$，“方程”
\[x^2=-xy-y+1-a,\qquad y^2=3y-2+3a\]
定义 $S$ 的自由且\Term{surjectivity}{满射}的\Term{group-extension}{扩张}；此时 $y^3+3xy^2+3x^2y-y=0$，前面的表达式给出 $p_{2\alpha+\beta}(a)\in H(S')$。\footnote{编者注：把 $p_{3\alpha+\beta}(a)$ 改正为 $p_{2\alpha+\beta}(a)$，并详述了论证的结尾。}于是已证明 $H$ 包含 $U_{2\alpha+\beta}$ 和 $U_{3\alpha+2\beta}$，且
\[E(x,y)=p_{\alpha+\beta}(xy)p_{3\alpha+\beta}(x^3y)\in H(S').\]
由于 $p_{\alpha+\beta}(xy)$ 与 $p_{3\alpha+\beta}$ 交换，便归结到前面的计算，即 $E(1,x)E(1,-x)=p_{3\alpha+\beta}(x^3-x)$，故 $U_{3\alpha+\beta}\subset H$，进而 $U_{\alpha+\beta}\subset H$。
\textbf{注记 6.2.1。} 上面的证明表明，可以把假设“$H$ 包含 $U_\alpha$ 和 $U_\beta$”替换为“$H$ 包含 $U_\alpha$ 或 $U_\beta$，且 $H$ 在 $\operatorname{int}(U_\alpha)$ 和 $\operatorname{int}(U_\beta)$ 下不变”。

\begin{SGAresult}{定理}{6.3}
设 $S$ 为\Term{scheme}{概形}，$G,G'$ 为两个\Term{semisimple}{半单} $S$-群，$B$（相应地，$B'$）为 $G$（相应地，$G'$）的 \Term{borel-subgroup}{Borel 子群}。任意\Term{isomorphism}{同构} $B\xrightarrow{\sim}B'$ 都由唯一\Term{isomorphism}{同构} $G\xrightarrow{\sim}G'$ 诱导。
\end{SGAresult}
断言对\Term{etale}{平展}\Term{topology}{拓扑}是\Term{local-scheme}{局部}的，可假设 $G$ \Term{split-group}{分裂}：$G=(G,T,M,R)$，$B$ 由 $R$ 的\Term{positive-root-system}{正根系} $R_+$ 定义。给定\Term{isomorphism}{同构} $u:B\xrightarrow{\sim}B'$ 诱导 $T$ 到 $B'$ 的一个\Term{maximal-torus}{极大环面} $T'$ 的\Term{isomorphism}{同构}，因而 $T'$ 也是 $G'$ 的\Term{maximal-torus}{极大环面}。给定\Term{isomorphism}{同构} $T\simeq D_S(M)$ 给出\Term{isomorphism}{同构} $T'\simeq D_S(M)$，其中 $R_+$ 的元素成为 $B'$ 相对于 $T'$ 的常根，从而 $R$ 的元素成为 $G'$ 相对于 $T'$ 的常根。由于 $G,G'$ \Term{semisimple}{半单}，\Term{coroot}{余根}由\Term{duality}{对偶性}确定（第 XXI 讲 1.2.5），这证明 $(T',M,R)$ 是 $G'$ 的一个\Term{split-group}{分裂化}，满足 $\mathscr R(G)=\mathscr R(G')$。
应用第 XXIII 讲 5.1（唯一性定理），得到唯一\Term{isomorphism}{同构} $G\xrightarrow{\sim}G'$，在 $T$ 及 $U_\alpha$（$\alpha\in\Delta$）上与 $u$ 相同。由 5.1，此\Term{isomorphism}{同构}在 $B$ 上的\Term{restriction}{限制}为 $u$。证毕。
\textbf{注记 6.3.1。} (i) 利用第 XXII 讲 4.1.9，并按该讲 4.2.12 的方式推理，可以在定理中把“\Term{isomorphism}{同构}”替换为“\Term{isogeny}{同源}”（相应地，“\Term{center}{中心}\Term{isogeny}{同源}”）。

(ii) 定理对\Term{reductive-group}{约化群}不成立。例如，取 $G=G'=\operatorname{SL}_2\times\mathbb G_m$，等同于下述矩阵群：
\[\left\{\begin{pmatrix}a&b&0\\c&d&0\\0&0&h\end{pmatrix}\;\middle|\;ad-bc=1,\ h\text{ 可逆}\right\};\]
\SourcePage{250}
取 $B=B'$ 为 $c=0$ 定义的 \Term{borel-subgroup}{Borel 子群}，$u$ 为如下 $B$ 的\Term{automorphism}{自同构}：$u(a,b,d,h)=(a,b,d,ah)$。
\begin{SGAresult}{推论}{6.4}
从由对 $(G,B)$（$G$ 为\Term{semisimple}{半单} $S$-群，$B$ 为其 \Term{borel-subgroup}{Borel 子群}）组成的\Term{category}{范畴}，到 $S$-群\Term{scheme}{概形}\Term{category}{范畴}的\Term{functor}{函子} $(G,B)\mapsto B$ 是\Term{fully-faithful}{全忠实}的（\Term{morphism}{态射}均取\Term{isomorphism}{同构}）。
\end{SGAresult}
\begin{SGAresult}{推论}{6.5}
设 $S$ 为\Term{scheme}{概形}，$G$ 为\Term{semisimple}{半单} $S$-群，$B$ 为 $G$ 的 \Term{borel-subgroup}{Borel 子群}，$B'$ 为对 (fpqc) \Term{local-isomorphism}{局部同构}于 $B$ 的 $S$-群。则 $B'$ 对\Term{locally-finite}{局部有限}\Term{etale}{平展}\Term{topology}{拓扑}\Term{local-isomorphism}{局部同构}于 $B$，\footnote{编者注：第 IV 讲 6.3 记为 (étf)。换言之，对每个 $s\in S$，存在 $s$ 的\Term{open-neighborhood}{开邻域} $U$ 及\Term{finite}{有限}\Term{etale}{平展}\Term{surjectivity}{满射} $V\to U$，使 $B'\times_SV\simeq B\times_SV$。}且存在\Term{semisimple}{半单} $S$-群 $G'$，使 $B'$ 为其 \Term{borel-subgroup}{Borel 子群}；$G'$ 在诱导 $B'$ 上\Term{identity-morphism}{恒等态射}的唯一\Term{isomorphism}{同构}意义下唯一。
\end{SGAresult}
\begin{SGAresult}{推论}{6.6}
设 $S$ 为\Term{scheme}{概形}，$G$ 为\Term{semisimple}{半单} $S$-群，$B$ 为 $G$ 的 \Term{borel-subgroup}{Borel 子群}。则 $\operatorname{Aut}_{S\text{-gr.}}(B)$ 可由\Term{affine}{仿射}\Term{smooth}{光滑} $S$-\Term{scheme}{概形}表示，$\operatorname{Autext}(B)$ 可由\Term{finite}{有限}\Term{etale}{平展} $S$-\Term{scheme}{概形}表示，且显然的\Term{morphism}{态射}诱导如下\Term{exact-sequence}{正合列}的\Term{isomorphism}{同构}：
\[
\begin{tikzcd}[column sep=small]
1\arrow[r]&B_{\mathrm{ad}}\arrow[r]\arrow[d,"\sim"']&\operatorname{Aut}_{S\text{-gr.}}(G,B)\arrow[r]\arrow[d,"\sim"']&\operatorname{Autext}(G)\arrow[r]\arrow[d,"\sim"']&1\\
1\arrow[r]&B_{\mathrm{ad}}\arrow[r]&\operatorname{Aut}_{S\text{-gr.}}(G)\arrow[r]&\operatorname{Autext}(B)\arrow[r]&1.
\end{tikzcd}
\]
\end{SGAresult}
这立即由 2.1 和上述结果得到。把第 1、2、3、4 节的结果在此框架中的类似结论展开，留给读者。
\textbf{注记 6.7。} 若 $S$ 为\Term{scheme}{概形}，$B$ 为 $S$-群，则 $B$ 至多是一个\Term{semisimple}{半单}群的 \Term{borel-subgroup}{Borel 子群}，而此\Term{semisimple}{半单}群由 $B$ 唯一确定。因此，尚须刻画那些确实是\Term{semisimple}{半单}群的 \Term{borel-subgroup}{Borel 子群}的 $S$-群 $B$。\footnote{编者注：可以问，若一个\Term{affine}{仿射}\Term{smooth}{光滑} $S$-群的每个\Term{geometric-fibre}{几何纤维}都是某个\Term{semisimple}{半单}群的 \Term{borel-subgroup}{Borel 子群}，它是否就是某个\Term{semisimple}{半单} $S$-群的 \Term{borel-subgroup}{Borel 子群}。（我们删去了原文给出的“反例”，其中 $S$ 是域 $k$ 上的\Term{duality}{对偶}数\Term{scheme}{概形}；M. Demazure 指出该反例有误。）}

\SGAsection{7}{当 $G$ 约化时函子 $\operatorname{Hom}_{S\text{-gr.}}(G,H)$ 的可表性}
\textbf{7.1. \Term{splitting}{分裂}情形}

\textbf{7.1.1。} 设 $S$ 为\Term{scheme}{概形}，$G$ 为带\Term{pinning}{钉扎}的 $S$-群，$T$ 为其\Term{maximal-torus}{极大环面}，$\Delta$ 为\Term{simple-root}{单根}\Term{set}{集合}；对 $\alpha\in\Delta$，设
\[u_\alpha\in U_\alpha^\times(S),\qquad w_\alpha\in\operatorname{Norm}_G(T)(S)\]
为\Term{pinning}{钉扎}所定义的元素。
\SourcePage{251}
另设 $H$ 为 $S$-群\Term{scheme}{概形}；我们关心\Term{functor}{函子} $\operatorname{Hom}_{S\text{-gr.}}(G,H)$，更确切地说，关心\Term{morphism}{态射}
\[q:\operatorname{Hom}_{S\text{-gr.}}(G,H)\longrightarrow\operatorname{Hom}_{S\text{-gr.}}(T,H).\]
因此，给定 $S$-\Term{group-morphism}{群态射} $f_T:T\to H$，考虑 $q^{-1}(f_T)=F$。它是如下定义的\Term{functor}{函子}：
\[F(S')=\{f\in\operatorname{Hom}_{S'\text{-gr.}}(G_{S'},H_{S'})\mid f=(f_T)_{S'}\text{ 于 }T_{S'}\}.\]
有一个 $S$-\Term{morphism-of-functors}{函子态射}
\[i:F\longrightarrow H^{2n},\]
其中 $n=\operatorname{Card}(\Delta)$，它在\Term{set}{集合}意义下由 $i(f)=(f(u_\alpha),f(w_\alpha))_{\alpha\in\Delta}$ 定义。由第 XXIII 讲 1.9，$i$ 还是\Term{monomorphism}{单态射}。
\begin{SGAresult}{命题}{7.1.2}
若 $H$ 在 $S$ 上\Term{separated}{分离}，则 $F$ \Term{representable}{可表}，且 $i$ 为\Term{closed-immersion}{闭浸入}。
\end{SGAresult}
相对\Term{representable}{可表}性的通常技巧\footnote{编者注：例如参见第 XXII 讲 5.8.1 的证明。}表明，只须证明：给定\Term{section}{截面} $v_\alpha,h_\alpha\in H(S)$（$\alpha\in\Delta$），使得存在延拓 $(f_T)_{S'}$ 且满足 $f(u_\alpha)=v_\alpha$、$f(w_\alpha)=h_\alpha$ 的 $S'$-\Term{homomorphism}{同态} $f:G_{S'}\to H_{S'}$ 的那些 $S$-\Term{scheme}{概形} $S'$，恰为通过 $S$ 的某个闭子\Term{scheme}{概形}\Term{resolution}{分解}的那些。显然可假设 $S$ \Term{affine}{仿射}。
为简化下文，称\Term{affine}{仿射}\Term{scheme}{概形}的\Term{morphism}{态射} $Y\to X$\footnote{编者注：把 $X\to Y$ 改正为 $Y\to X$。}满足条件 (L)，若 $Y$ 为一个 $\mathscr O(X)$-\Term{algebra}{代数}的谱，且此\Term{algebra}{代数}是自由 $\mathscr O(X)$-模。显然，在\Term{affine}{仿射}\Term{scheme}{概形}\Term{category}{范畴}中，满足 (L) 的\Term{morphism}{态射}的乘积或复合仍满足 (L)，且 (L) 在底\Term{group-extension}{扩张}下\Term{stable}{稳定}。
\begin{SGAresult}{引理}{7.1.3}
设 $S$ \Term{affine}{仿射}，$\alpha\in\Delta$。考虑在\Term{set}{集合}意义下由 $a(t,t')=\alpha(t)+\alpha(t')$ 定义的\Term{morphism}{态射}
\[a:T\times_ST\longrightarrow\mathbb G_{a,S}.\]
(i) $a$ \Term{faithfully-flat}{忠实平坦}且\Term{finite-presentation}{有限表示}。

(ii) 设 $R$ 为 $a$ 的\Term{fiber}{纤维}平方。\Term{structural-morphism}{结构态射} $R\to S$ 满足 (L)。
\end{SGAresult}
\SourcePage{252}
首先，$a$ 显然\Term{surjectivity}{满射}，因为 $\alpha:T\to\mathbb G_{m,S}$ \Term{surjectivity}{满射}。$a$ 显然\Term{finite-presentation}{有限表示}。由于 $\alpha$ 满足 (L)，要证明 (i)、(ii)，只须证明在\Term{set}{集合}意义下由 $u(x,y)=x+y$ 定义的\Term{morphism}{态射}
\[u:\mathbb G_{m,S}^2\longrightarrow\mathbb G_{a,S}\]
满足 (L)。换言之，只须验证对任意环 $A$，环 $A[X,Y,X^{-1},Y^{-1}]$ 是其子环 $A[X+Y]$ 上的\Term{free-module}{自由模}。\footnote{编者注：简化了原文的证明。}而 $A[X,Y]$ 在子环 $B=A[X+Y,XY]$ 上是以 $\{1,X\}$ 为基的\Term{free-module}{自由模}（因为 $X^2-(X+Y)X+XY=0$），所以 $A[X^{\pm1},Y^{\pm1}]=A[X,Y]_{(XY)}$ 在 $B_{XY}=A[X+Y,(XY)^{\pm1}]$ 上自由，基为 $\{1,X\}$；它在 $A[X+Y]$ 上也自由，基为 $(XY)^p$ 和 $(XY)^pX$（$p\in\mathbb Z$）。
\begin{SGAresult}{引理}{7.1.4}
设 $\alpha\in\Delta$，$b:T\times_ST\to H$ 为在\Term{set}{集合}意义下由
\[b(t,t')=\bigl(\operatorname{int}(f_T(t))v_\alpha\bigr)\bigl(\operatorname{int}(f_T(t'))v_\alpha\bigr)\]
定义的\Term{morphism}{态射}。设 $f_\alpha:U_\alpha\to H$ 为 $S$-\Term{morphism}{态射}。以下条件\Term{equivalence}{等价}：\footnote{编者注：以下 $U_\alpha$ 的群律用加法记号；即若 $p_\alpha:\mathbb G_{a,S}\xrightarrow{\sim}U_\alpha$ 为满足 $p_\alpha(1)=u_\alpha$ 的\Term{isomorphism}{同构}，且 $x=p_\alpha(z)$，则 $\alpha(t)x$（相应地，$a(t,t')u_\alpha$）表示 $p_\alpha(\alpha(t)z)$（相应地，$p_\alpha(a(t,t'))=\alpha(t)u_\alpha+\alpha(t')u_\alpha$）。}

(i) $f_\alpha$ 为\Term{group-morphism}{群态射}；有 $f_\alpha(u_\alpha)=v_\alpha$，且对所有 $t\in T(S')$、$x\in U_\alpha(S')$、$S'\to S$，有
\[f_T(t)f_\alpha(x)f_T(t)^{-1}=f_\alpha(\alpha(t)x).\]
(ii) 有 $f_\alpha(u_\alpha)=v_\alpha$，且对所有 $t,t'\in T(S')$、$S'\to S$，有关系
\[f_\alpha(a(t,t')u_\alpha)=b(t,t').\]
\end{SGAresult}
若 $f_\alpha$ 满足 (i)，则有 $f_\alpha(\alpha(t)u_\alpha)=\operatorname{int}(f_T(t))v_\alpha$，立即推出 (ii)。反之，设 (ii) 成立，证明 (i) 的各项条件；先证最后一条。由于 $a$ 是 (fpqc) \Term{covering}{覆盖}，只须证明对 $t,t',t''\in T(S')$，有
\[f_T(t)f_\alpha(a(t',t'')u_\alpha)f_T(t)^{-1}=f_\alpha(\alpha(t)a(t',t'')u_\alpha);\]
这也可写成 $f_T(t)b(t',t'')f_T(t)^{-1}=b(tt',tt'')$，由 $b$ 的定义显然成立。还须证明 $f_\alpha$ 为\Term{group-morphism}{群态射}。上述性质立即给出
\[\begin{aligned}
f_\alpha(\alpha(t)u_\alpha)f_\alpha(\alpha(t')u_\alpha)
&=(f_T(t)f_\alpha(u_\alpha)f_T(t)^{-1})\cdot(f_T(t')f_\alpha(u_\alpha)f_T(t')^{-1})\\
&=b(t,t')=f_\alpha(\alpha(t)u_\alpha+\alpha(t)\alpha(t')u_\alpha).
\end{aligned}\]
于是只要 $x,x'$ 是 $U_\alpha$ 的\Term{open-subset}{开集} $U_\alpha^\times$ 的\Term{section}{截面}，就有 $f_\alpha(x+x')=f_\alpha(x)f_\alpha(x')$；该\Term{open-subset}{开集}在\Term{scheme}{概形}意义下\Term{dense}{稠密}，故由第 XVIII 讲 1.4 得到结论。

\SourcePage{253}
\textbf{7.1.5。} 暂时固定 $\alpha\in\Delta$。\Term{morphism}{态射} $a$ \Term{faithfully-flat}{忠实平坦}且\Term{quasi-compact}{拟紧}，因此 $\mathbb G_{a,S}$ 等同于 $T\times_ST$ 对 $a$ 所定义的\Term{equivalence-relation}{等价关系} $R$ 的商。把此\Term{equivalence-relation}{等价关系}写为
\[\begin{tikzcd}R\arrow[r,shift left,"i_1"]\arrow[r,shift right,"i_2"']&T\times_ST.\end{tikzcd}\]
\Term{morphism}{态射} $b$ 通过 $a$ \Term{resolution}{分解}\footnote{编者注：即 7.1.4 的条件 (ii) 成立。}的充要条件是 $b\circ i_1=b\circ i_2$；换言之，若 $K$ 表示\Term{morphism}{态射}对
\[\begin{tikzcd}R\arrow[r,shift left,"b\circ i_1"]\arrow[r,shift right,"b\circ i_2"']&H\end{tikzcd}\]
的核，则充要条件是 $K=R$。而假设 $H$ 在 $S$ 上\Term{separated}{分离}，故 $K$ 是 $R$ 的闭子\Term{scheme}{概形}。另一方面，回顾 $R$ 在 $S$ 上\Term{essentially-free}{本质自由}（7.1.3）。

\textbf{7.1.6。} 由上可知，若 $S'$ 为 $S$-\Term{scheme}{概形}，则在 $S'$ 上存在满足 7.1.4 (i) 的 $f_\alpha$（此时必然唯一）的充要条件是 $K_{S'}=R_{S'}$，且所得\Term{morphism}{态射} $f_\alpha$ 满足 $f_\alpha(u_\alpha)=v_\alpha$。
第一个条件\Term{equivalence}{等价}于 $S'\to S$ 通过 $S$ 的某个闭子\Term{scheme}{概形}\Term{resolution}{分解}（第 VIII 讲 6.4\footnote{编者注：另见第 VI$_{\mathrm B}$ 讲 6.2.3 的增补。}）；若把 $S$ 换成此闭子\Term{scheme}{概形}，则第二个条件又定义 $S$ 的一个闭子\Term{scheme}{概形}（$H$ 的两个\Term{section}{截面}相等，而 $H$ 在 $S$ 上\Term{separated}{分离}）。
因而把 $S$ 换成此闭子\Term{scheme}{概形}后，可假设存在满足 7.1.4 (i) 的\Term{morphism}{态射} $f_\alpha:U_\alpha\to H$。取各个 $\alpha\in\Delta$ 所得 $S$ 的子\Term{scheme}{概形}之交，即可假设此条件对所有 $\alpha\in\Delta$ 都成立。

\textbf{7.1.7。} 同样，对每个 $\alpha\in\Delta$ 考虑两个 $S$-\Term{group-morphism}{群态射}
\[f_T\circ\operatorname{int}(w_\alpha),\quad\operatorname{int}(h_\alpha)\circ f_T:T\longrightarrow H.\]
由于 $H$ 在 $S$ 上\Term{separated}{分离}，且 $T$ 在 $S$ 上\Term{essentially-free}{本质自由}，与前面相同的推理表明，把 $S$ 换成一个闭子\Term{scheme}{概形}后，可假设对每个 $\alpha\in\Delta$ 都有
\[f_T\circ\operatorname{int}(w_\alpha)=\operatorname{int}(h_\alpha)\circ f_T.\]

\textbf{7.1.8。} 现在利用\Term{generator}{生成元}与关系定理（第 XXIII 讲 3.5），可见存在满足所需条件的\Term{group-homomorphism}{群同态} $f:G\to H$ 的充要条件是：$H$ 的\Term{section}{截面} $h_\alpha,v_\alpha,f_T(\alpha^*(-1))$（$\alpha\in\Delta$）满足某个\Term{finite}{有限}的\Term{algebra}{代数}关系集：
\[R_i\bigl((h_\alpha)_\alpha,(v_\alpha)_\alpha,(f_T(\alpha^*(-1)))_\alpha\bigr)=e,\qquad i=1,\ldots,n.\]
\SourcePage{254}
由于 $H$ 在 $S$ 上\Term{separated}{分离}，这又定义 $S$ 的一个闭子\Term{scheme}{概形}，证明完成。
\begin{SGAresult}{推论}{7.1.9}
设 $S$ 为\Term{scheme}{概形}，$G$ 为\Term{split-group}{分裂}\Term{reductive-group}{约化} $S$-群，$T$ 为其\Term{maximal-torus}{极大环面}，$H$ 为在 $S$ 上\Term{separated}{分离}的 $S$-群\Term{scheme}{概形}。另设 $P$ 为\Term{morphism}{态射}的一个性质，使：

(i) \Term{closed-immersion}{闭浸入}满足 $P$；

(ii) 两个满足 $P$ 的\Term{morphism}{态射}的复合也满足 $P$；

(iii) $P$ 在\Term{base-change}{基变换}下\Term{stable}{稳定}；

(iv) \Term{structural-morphism}{结构态射} $H\to S$ 满足 $P$。

则\Term{canonical}{典范}\Term{morphism}{态射}
\[\operatorname{Hom}_{S\text{-gr.}}(G,H)\longrightarrow\operatorname{Hom}_{S\text{-gr.}}(T,H)\]
相对地由满足 $P$ 的\Term{separated}{分离}\Term{morphism}{态射}表示。\footnote{编者注：即对任意 $S'\to\operatorname{Hom}_{S\text{-gr.}}(T,H)$，其中 $S'$ \Term{representable}{可表}，$\operatorname{Hom}_{S\text{-gr.}}(G,H)\times_SS'$ \Term{representable}{可表}，且 $S$-\Term{scheme}{概形}\Term{morphism}{态射} $\operatorname{Hom}_{S\text{-gr.}}(G,H)\times_SS'\to S'$ \Term{separated}{分离}且满足 $P$。}
\end{SGAresult}
事实上，可以假设 $G$ 带\Term{pinning}{钉扎}；\Term{structural-morphism}{结构态射} $H^{2n}\to S$ 满足 $P$，由 7.1.2 即得结论。
\begin{SGAresult}{推论}{7.1.10}
设 $S$ 为\Term{scheme}{概形}，$G$ 为\Term{split-group}{分裂}\Term{reductive-group}{约化} $S$-群，$H$ 为\Term{smooth}{光滑}且\Term{fiber}{纤维}\Term{affine}{仿射}的 $S$-群\Term{scheme}{概形}。\footnote{编者注：删去了多余的 $H$ \Term{quasi-separated}{拟分离}的假设。}则\Term{functor}{函子} $\operatorname{Hom}_{S\text{-gr.}}(G,H)$ 可由在 $S$ 上\Term{separated}{分离}且\Term{locally-finite-presentation}{局部有限表示}的 $S$-\Term{scheme}{概形}表示。
\end{SGAresult}
事实上，由于 $H$ 在 $S$ 上\Term{smooth}{光滑}，可考虑其\Term{identity-component}{中性分支} $H^0$；它在 $S$ 上\Term{fiber}{纤维}\Term{affine}{仿射}、\Term{smooth}{光滑}、\Term{separated}{分离}且\Term{finite-presentation}{有限表示}（第 VI$_{\mathrm B}$ 讲 3.10 和 5.5）。由于 $G$ 的\Term{fiber}{纤维}\Term{connected}{连通}，可以把 $H$ 换成 $H^0$。此时 $G,H$ 均\Term{finite-presentation}{有限表示}，可假设 $S$ \Term{noetherian}{诺特}，并在 7.1.9 中取 $P$ 为“\Term{finite-type}{有限型}”。而由第 XV 讲 8.9 已知，$\operatorname{Hom}_{S\text{-gr.}}(T,H)$ 可由\Term{separated}{分离}且\Term{locally-finite-type}{局部有限型}的 $S$-\Term{scheme}{概形}表示。
\textbf{注记 7.1.11。} 若 $H$ 在 $S$ 上\Term{affine}{仿射}，则可把对第 XV 讲的引用换成第 XI 讲 4.2。

\textbf{7.2. 一般情形}
\begin{SGAresult}{命题}{7.2.1}
设 $S$ 为\Term{scheme}{概形}，$G$ 为\Term{reductive-group}{约化} $S$-群，$T$ 为 $G$ 的\Term{maximal-torus}{极大环面}。另设 $H$ 为 $S$-群\Term{scheme}{概形}，其\Term{structural-morphism}{结构态射} $H\to S$ 满足条件：

(+) 每个点 $s\in S$ 都有\Term{open-neighborhood}{开邻域} $U$，使\Term{morphism}{态射} $H_U\to U$ \Term{quasi-projective}{拟射影}。

则\Term{canonical}{典范}\Term{morphism}{态射}
\[\operatorname{Hom}_{S\text{-gr.}}(G,H)\longrightarrow\operatorname{Hom}_{S\text{-gr.}}(T,H)\]
相对地由满足 (+) 的\Term{morphism}{态射}表示。
\end{SGAresult}
\SourcePage{255}
当 $G$ 相对于 $T$ \Term{split-group}{可分裂}时，在 7.1.9 中取 $P$ 为上述性质 (+)。当 $G$ \Term{locally-isotrivial}{局部等平凡}，例如\Term{semisimple}{半单}时（4.2.2），注意到定理的断言对于\Term{locally-finite}{局部有限}\Term{etale}{平展}\Term{topology}{拓扑}是\Term{local-scheme}{局部}的（事实上，\Term{quasi-projective}{拟射影}性对于整体\Term{finite}{有限}\Term{etale}{平展}\Term{topology}{拓扑}是\Term{local-scheme}{局部}的，并保证此\Term{topology}{拓扑}下\Term{descent}{下降}的\Term{effective}{有效}性，参见 SGA 1，第 VIII 讲 7.7）。最后，在一般情形使用下述引理：
\begin{SGAresult}{引理}{7.2.2}
设 $S$ 为\Term{scheme}{概形}，$G$ 为\Term{reductive-group}{约化} $S$-群，$T$ 为 $G$ 的\Term{maximal-torus}{极大环面}，$G'$ 为 $G$ 的\Term{derived-group}{导群}（第 XXII 讲 6.2），$T'=T\cap G$ 为 $G'$ 的对应于 $T$ 的\Term{maximal-torus}{极大环面}（第 XXII 讲 6.2.8）。则图
\[\begin{tikzcd}G&G'\arrow[l]\\T\arrow[u]&T'\arrow[l]\arrow[u]\end{tikzcd}\]
在 $S$-群\Term{sheaf}{层}\Term{category}{范畴}中为\Term{amalgamated-sum}{融合和}；即对任意 $S$-群\Term{sheaf}{层} $H$，以下方图为\Term{cartesian-diagram}{笛卡尔图}：
\[\begin{tikzcd}
\operatorname{Hom}_{S\text{-gr.}}(G,H)\arrow[r]\arrow[d]&\operatorname{Hom}_{S\text{-gr.}}(G',H)\arrow[d]\\
\operatorname{Hom}_{S\text{-gr.}}(T,H)\arrow[r]&\operatorname{Hom}_{S\text{-gr.}}(T',H).
\end{tikzcd}\]
\end{SGAresult}
事实上，若 $\operatorname{rad}(G)$ 为 $G$ 的\Term{jacobson-radical}{根基}，则 $\operatorname{rad}(G)\subset T$，所以
\[\operatorname{rad}(G)\cap G'=\operatorname{rad}(G)\cap T'=K,\]
且 $G$ 中的乘法诱导\Term{isogeny}{同源}（第 XXII 讲 6.2）
\[i:G'\times_S\operatorname{rad}(G)\longrightarrow G,\qquad j:T'\times_S\operatorname{rad}(G)\longrightarrow T.\]
设 $f_{G'}:G'\to H$、$f_T:T\to H$、$f_{T'}:T'\to H$ 为 $S$-\Term{group-morphism}{群态射}，满足 $f_{G'}|_{T'}=f_T|_{T'}=f_{T'}$。证明存在唯一 $S$-\Term{group-morphism}{群态射} $f:G\to H$ 诱导 $f_{G'}$ 和 $f_T$。置 $f_{\mathrm{rad}}=f_T|_{\operatorname{rad}(G)}$，\footnote{编者注：把 $f_r$ 换成了 $f_{\mathrm{rad}}$。}并置
\[f_1=f_{G'}\cdot f_{\mathrm{rad}}:G'\times_S\operatorname{rad}(G)\longrightarrow H.\]
$f$ 存在（而且显然唯一）的充要条件是 $f_1$ 在 $i$ 的核上诱导恒等元，即 $f_{G'}$ 与 $f_{\mathrm{rad}}$ 在 $K$ 上相同；但用同样的推理，$f_T$ 的存在表明 $f_{T'}$ 与 $f_{\mathrm{rad}}$ 在 $K$ 上相同。只须再注意到 $f_{G'}$ 与 $f_{T'}$ 在 $K\subset T'$ 上显然相同。
现按 7.1.10 的方式推理，由 7.2.1 得到：
\begin{SGAresult}{推论}{7.2.3}
设 $S$ 为\Term{scheme}{概形}，$G$ 为\Term{reductive-group}{约化} $S$-群，$H$ 为在 $S$ 上\Term{smooth}{光滑}、\Term{quasi-projective}{拟射影}且\Term{fiber}{纤维}\Term{affine}{仿射}的 $S$-群\Term{scheme}{概形}。则 $\operatorname{Hom}_{S\text{-gr.}}(G,H)$ 可由在 $S$ 上\Term{separated}{分离}且\Term{locally-finite-presentation}{局部有限表示}的 $S$-\Term{scheme}{概形}表示。
\end{SGAresult}
\SourcePage{256}
\textbf{7.3. \Term{characteristic-zero}{特征零}所特有的现象}

若 $G,H$ 为两个\Term{smooth}{光滑} $S$-群，$\mathfrak g,\mathfrak h$ 为其 \Term{lie-algebra}{Lie 代数}，则有\Term{canonical}{典范}\Term{morphism}{态射}
\[\operatorname{Lie}:\operatorname{Hom}_{S\text{-gr.}}(G,H)\longrightarrow\operatorname{Hom}_{\mathscr O_S\text{-Lie}}(\mathfrak g,\mathfrak h),\]
右端的定义显然。
\begin{SGAresult}{命题}{7.3.1}
设 $S$ 为\Term{characteristic-zero}{特征零}\Term{scheme}{概形}（即 $\mathbb Q$-\Term{scheme}{概形}），$G$ 为\Term{reductive-group}{约化} $S$-群，$H$ 为在 $S$ 上\Term{smooth}{光滑}、\Term{quasi-projective}{拟射影}且\Term{fiber}{纤维}\Term{affine}{仿射}的 $S$-群\Term{scheme}{概形}。

(i) $\operatorname{Hom}_{S\text{-gr.}}(G,H)$ 可由在 $S$ 上\Term{smooth}{光滑}且\Term{separated}{分离}的 $S$-\Term{scheme}{概形}表示。

(ii) 若 $G$ \Term{semisimple}{半单}，则该\Term{scheme}{概形}在 $S$ 上\Term{affine}{仿射}且\Term{finite-presentation}{有限表示}。

(iii) 若 $G$ \Term{simply-connected}{单连通}（第 XXII 讲 4.3.3），则\Term{canonical}{典范}\Term{morphism}{态射}
\[\operatorname{Lie}:\operatorname{Hom}_{S\text{-gr.}}(G,H)\longrightarrow\operatorname{Hom}_{\mathscr O_S\text{-Lie}}(\mathfrak g,\mathfrak h)\]
为\Term{isomorphism}{同构}。
\end{SGAresult}
\begin{SGAresult}{引理}{7.3.2}
设 $k$ 为\Term{characteristic-zero}{特征零}域，$G$ 为\Term{reductive-group}{约化} $k$-群，$V$ 为\Term{finite}{有限}维 $k$-\Term{vector-space}{向量空间}，$G\to\operatorname{GL}(V)$ 为 $G$ 在 $V$ 中的\Term{linear-representation}{线性表示}。则
\[H^0(G,V)=H^0(\mathfrak g,V),\qquad H^i(G,V)=0\quad(i>0).\]
\end{SGAresult}
第一个等式一般地对\Term{smooth}{光滑}\Term{connected}{连通}群成立；\footnote{编者注：显然有 $V^G\subset V^{\mathfrak g}$。另一方面，$H=\operatorname{Centr}_G(V^{\mathfrak g})$ 是 $G$ 的\Term{closed-subgroup}{闭子群}，因 $\operatorname{car}(k)=0$ 而\Term{smooth}{光滑}；由第 II 讲 5.3.1，有 $\operatorname{Lie}(H)=\operatorname{Centr}_{\mathfrak g}(V^{\mathfrak g})=\mathfrak g$，由于 $H$ \Term{smooth}{光滑}且 $G$ \Term{connected}{连通}，这推出 $H=G$，所以 $V^{\mathfrak g}\subset V^G$（另见 [DG70]，\S II.6，命题 2.1 (c)）。}在\Term{reductive-group}{约化群}情形，可如下证明：可假设 $k$ \Term{algebraically-closed}{代数闭}，从而 $G$ \Term{split-group}{可分裂}，故 $G$ 由\Term{isomorphism}{同构}于 $\mathbb G_{m,k}$ 的\Term{subgroup}{子群}生成，\footnote{编者注：因为 $G$ 的\Term{maximal-torus}{极大环面}的并在 $G$ 中\Term{dense}{稠密}，参见 Bible，\S 6.5，定理 5（即 [Ch05]，\S 6.6，定理 6）。}因而只须对这个群验证断言，这是容易的。
由第一个等式可知，当 $G$ \Term{semisimple}{半单}时，$H^0(G,V)$ 是关于 $V$ 的\Term{exact-functor}{正合函子}；因为此时 $\mathfrak g$ 是\Term{semisimple}{半单} \Term{lie-algebra}{Lie 代数}，应用 [BLie]，\S I.6，习题 1 (b) 即可。$G$ \Term{reductive-group}{约化}时，断言仍成立：引入 $G$ 的\Term{jacobson-radical}{根基} $R$\footnote{编者注：注意 $R$ 是一个\Term{torus}{环面}。}和\Term{semisimple}{半单}\Term{quotient-group}{商群} $G/R$，有 $H^0(G,V)=H^0(G/R,H^0(R,V))$，而 $H^0(G,-)$ 是两个\Term{exact-functor}{正合函子}的复合。再应用第 I 讲 5.3.1，得到 $H^i(G,V)=0$（$i>0$）。
\textbf{注记 7.3.3。} 在上述条件下，若 $G$ \Term{semisimple}{半单}，则 $H^1(\mathfrak g,V)=H^2(\mathfrak g,V)=0$；参见 Bourbaki，上引处 (b)、(d)。

\SourcePage{257}
\textbf{7.3.4。} 现在证明命题。由 7.2.3 已知，$\operatorname{Hom}_{S\text{-gr.}}(G,H)$ 可由在 $S$ 上\Term{separated}{分离}且\Term{locally-finite-presentation}{局部有限表示}的 $S$-\Term{scheme}{概形}表示；要证明其\Term{smooth}{光滑}，只须证明其\Term{infinitesimal}{无穷小}\Term{smooth}{光滑}（第 XI 讲 1.8），\footnote{编者注：即对每个\Term{closed-point}{闭点}为 $s'$ 的\Term{locally-artinian}{局部 Artin} $S$-\Term{scheme}{概形} $S'$，任意 $\kappa(s')$-\Term{group-morphism}{群态射} $G_{s'}\to H_{s'}$ 都可\Term{lifting}{提升}为 $S'$-\Term{group-morphism}{群态射} $G_{S'}\to H_{S'}$。由第 III 讲 2.8，这由 $H^2(G_{s'},V)$ 的消失推出，其中 $V=\operatorname{Lie}(H_{s'}/\kappa(s'))$。}这由 7.3.2 及第 III 讲 2.8 (i) 得到。因而 (i) 已证。
证明 (ii) 可由 (iii) 推出。首先只须证明 $\operatorname{Hom}_{S\text{-gr.}}(G,H)$ 在 $S$ 上\Term{affine}{仿射}；由 (i)，它在 $S$ 上\Term{smooth}{光滑}，故将\Term{finite-presentation}{有限表示}。无论如何，$\operatorname{Hom}_{\mathscr O_S\text{-Lie}}(\mathfrak g,\mathfrak h)$ 可由
\[\operatorname{Hom}_{\mathscr O_S\text{-mod.}}(\mathfrak g,\mathfrak h)\simeq\mathbf W(\mathfrak g^\vee\otimes_{\mathscr O_S}\mathfrak h)\]
的一个闭子\Term{scheme}{概形}表示，后者是\Term{affine}{仿射} $S$-\Term{scheme}{概形}，所以当 $G$ \Term{simply-connected}{单连通}时，所需结论立即出现。
一般情形，可假设 $G$ \Term{split-group}{分裂}，所以 $G\simeq G_S^{\mathrm{Ép}}(\mathscr R)$；引入\Term{simply-connected}{单连通}\Term{root-datum}{根数据} $\operatorname{sc}(\mathscr R)$（第 XXI 讲 6.5.5），利用存在性定理（第 XXV 讲 1.1），构造一个\Term{simply-connected}{单连通} $S$-群 $\overline G$ 和\Term{center}{中心}\Term{isogeny}{同源} $\overline G\to G$。此\Term{isogeny}{同源}的核 $K$ 是\Term{finite}{有限}\Term{diagonalizable}{可对角化} $S$-群（由于 $S$ \Term{characteristic-zero}{特征零}，故为\Term{twisted-constant}{扭曲常值} $S$-群），且 $\operatorname{Hom}_{S\text{-gr.}}(K,H)$（显然）可由\Term{separated}{分离} $S$-\Term{scheme}{概形}表示（若 $K\simeq(\mathbb Z/n\mathbb Z)_S$，则 $\operatorname{Hom}_{S\text{-gr.}}(K,H)$ \Term{isomorphism}{同构}于 $\operatorname{Ker}(H\xrightarrow{n}H)$）。有“带点 $S$-\Term{scheme}{概形}”的\Term{exact-sequence}{正合列}
\[1\longrightarrow\operatorname{Hom}_{S\text{-gr.}}(G,H)\xrightarrow{u}\operatorname{Hom}_{S\text{-gr.}}(\overline G,H)\longrightarrow\operatorname{Hom}_{S\text{-gr.}}(K,H),\]
所以 $u$ 是\Term{closed-immersion}{闭浸入}，从而 $\operatorname{Hom}_{S\text{-gr.}}(G,H)$ 在 $S$ 上\Term{affine}{仿射}。

\textbf{7.3.5。} 最后证明 (iii)。按 (i) 的证明推理并利用 7.3.3，可以证明 $S$-\Term{scheme}{概形} $\operatorname{Hom}_{\mathscr O_S\text{-Lie}}(\mathfrak g,\mathfrak h)$ 在 $S$ 上\Term{smooth}{光滑}。因此为证明 (iii)，可假设 $S=\operatorname{Spec}(k)$，其中 $k$ 为\Term{characteristic-zero}{特征零}\Term{algebraically-closed-field}{代数闭域}；甚至只须证明\Term{morphism}{态射} $\operatorname{Lie}$ 在 $k$ 值点上\Term{bijection}{双射}，且在两个对应点处的\Term{tangent-space}{切空间}间诱导\Term{bijection}{双射}。先证明
\[\operatorname{Hom}_{k\text{-gr.}}(G,H)\longrightarrow\operatorname{Hom}_{k\text{-Lie}}(\mathfrak g,\mathfrak h)\]
为\Term{bijection}{双射}。
若 $u:G\to H$ 为 $k$-\Term{group-morphism}{群态射}，其图像为 $G\times_kH$ 的\Term{connected}{连通}\Term{subgroup}{子群}，由其 \Term{lie-algebra}{Lie 代数}确定，而此 \Term{lie-algebra}{Lie 代数}是 $\operatorname{Lie}(u)$ 的图像；故映射\Term{injectivity}{单射}。反之，若 $v:\mathfrak g\to\mathfrak h$ 为 \Term{lie-algebra}{Lie 代数}\Term{morphism}{态射}，$v$ 的图像 $\mathfrak k$ 为 $\mathfrak g\times\mathfrak h$ 的一个 \Term{lie-subalgebra}{Lie 子代数}，\Term{isomorphism}{同构}于 $\mathfrak g$。特别地，因为 $\mathfrak g$ 等于其自身的导\Term{algebra}{代数}，$\mathfrak k$ 也如此；由 Chevalley 的一个定理（[Ch51]，\S 14，定理 15），$\mathfrak k$ 是 $G\times_kH$ 的某个\Term{connected}{连通}\Term{subgroup}{子群} $K$ 的 \Term{lie-algebra}{Lie 代数}。
\SourcePage{258}
\footnote{编者注：另见 [Bo91]，II 7.9。另外增添了随后一句。}又因为 $\mathfrak k\simeq\mathfrak g$ \Term{semisimple}{半单}，故 $G$ 是\Term{semisimple}{半单} $k$-群。由于
\[\dim(K)=\dim(\mathfrak k)=\dim(\mathfrak g)=\dim(G),\]
且 $K\cap H$ \Term{finite}{有限}（因其 \Term{lie-algebra}{Lie 代数}为零），\Term{canonical}{典范}\Term{morphism}{态射} $\operatorname{pr}_1:K\to G$ \Term{finite}{有限}且\Term{dominant}{占优}；由于 $G$ \Term{connected}{连通}，$\operatorname{pr}_1$ \Term{surjectivity}{满射}，故为\Term{isogeny}{同源}。因为 $G$ \Term{simply-connected}{单连通}而 $K$ \Term{semisimple}{半单}，由第 XXI 讲 6.2.7，$\operatorname{pr}_1$ 为\Term{isomorphism}{同构}，故 $K$ 为一个 $k$-\Term{group-morphism}{群态射} $u:G\to H$ 的图像，且 $\operatorname{Lie}(u)=v$。
最后，设 $u:G\to H$ 为 $k$-\Term{group-morphism}{群态射}。$\operatorname{Hom}_{k\text{-gr.}}(G,H)$ 在 $u$ 处的\Term{tangent-space}{切空间}等同于 $Z^1(G,\mathfrak h)$（参见第 II 讲 4.2；$G$ 通过 $\operatorname{Ad}_H\circ u$ \Term{action}{作用}于 $\mathfrak h$）；同样可证明，$\operatorname{Hom}_{k\text{-Lie}}(\mathfrak g,\mathfrak h)$ 在 $\operatorname{Lie}(u)$ 处的\Term{tangent-space}{切空间}等同于 $Z^1(\mathfrak g,\mathfrak h)$。因此须证明\Term{canonical}{典范}映射 $Z^1(G,\mathfrak h)\to Z^1(\mathfrak g,\mathfrak h)$ \Term{bijection}{双射}。而有\Term{commutative-diagram}{交换图}
\[\begin{tikzcd}[column sep=small]
0\arrow[r]&\mathfrak h^G\arrow[r]\arrow[d]&\mathfrak h\arrow[r]\arrow[d,equal]&Z^1(G,\mathfrak h)\arrow[r]\arrow[d]&H^1(G,\mathfrak h)\\
0\arrow[r]&\mathfrak h^{\mathfrak g}\arrow[r]&\mathfrak h\arrow[r]&Z^1(G,\mathfrak h)\arrow[r]&H^1(\mathfrak g,\mathfrak h).
\end{tikzcd}\]
由 7.3.2 和 7.3.3，$\mathfrak h^G=\mathfrak h^{\mathfrak g}$ 且 $H^1(G,\mathfrak h)=0=H^1(\mathfrak g,\mathfrak h)$，从而得到所需结论。
\textbf{注记 7.3.6。} 很可能，若 $S$ 为\Term{locally-noetherian}{局部诺特}\Term{scheme}{概形}，$G$ 为\Term{semisimple}{半单}\Term{simply-connected}{单连通} $S$-群，$H$ 为\Term{smooth}{光滑} $S$-群\Term{scheme}{概形}，$\widehat G,\widehat H$ 为它们沿\Term{unit-section}{单位截面}的\Term{formal-completion}{形式完备化}，则任意\Term{homomorphism}{同态} $\widehat G\to\widehat H$ 都来自唯一的 $G$ 到 $H$ 的\Term{homomorphism}{同态}；这将推广 7.3.1 (iii)。当 $S$ 是一个域 $k$ 的谱且 $H$ \Term{affine}{仿射}\Term{finite-type}{有限型}时，这由 Dieudonné 的一个定理得到（[Di57]，\S 15，定理 4）。\footnote{编者注：\Term{form}{形式}群 $\widehat G,\widehat H$ 的定义见第 VII$_{\mathrm B}$ 讲。设 $S=\operatorname{Spec}(k)$，其中 $k$ 为域。由该讲 2.2.1，给定\Term{form}{形式} $k$-\Term{group-morphism}{群态射} $v:\widehat G\to\widehat H$，\Term{equivalence}{等价}于给定 Hopf $k$-\Term{algebra}{代数}\Term{morphism}{态射} $\phi:\operatorname{Dist}(G)\to\operatorname{Dist}(H)$，其中 $\operatorname{Dist}(G)$ 表示 $G$ 的（原点处）分布\Term{algebra}{代数}，参见第 VII$_{\mathrm A}$ 讲 2.1，或 [DG70]，\S II.4，6.1，或 [Ja87]，I 7.7（[Di57] 和 [Ta75] 中称为 $G$ 的“超\Term{algebra}{代数}”）。[Di57] 的定理 4（另见 [Ta75]，0.3.4 (f)、(g)）在此背景下推广了 7.3.5 所用的 Chevalley 定理；由此得到，存在 $G\times H$ 的一个\Term{connected}{连通}闭 $k$-\Term{subgroup}{子群} $K$，使 $\widehat K$ 等于 $v$ 的图像；由于 $\operatorname{Dist}(K)\to\operatorname{Dist}(G)$ 为\Term{isomorphism}{同构}，$K\to G$ 为\Term{finite}{有限}\Term{etale}{平展}\Term{morphism}{态射}。由此推出 $K$ \Term{semisimple}{半单}，再由第 XXI 讲 6.2.7，$K\to G$ 为\Term{isomorphism}{同构}（因 $G$ \Term{simply-connected}{单连通}）；另见 [Ta75]，1.8 和 2.2。更一般地，该文研究满足性质 (SC) 的 $k$-群 $G$：任意\Term{finite}{有限}\Term{etale}{平展} $k$-\Term{group-morphism}{群态射} $G'\to G$，若 $G'$ \Term{connected}{连通}，则是\Term{isomorphism}{同构}。最后，上述结果表明，任意\Term{finite}{有限}维 $\operatorname{Dist}(G)$-模 $V$ 都以唯一方式成为 $G$-模；对\Term{split-group}{分裂}\Term{reductive-group}{约化} $k$-群 $G$（甚至 $G$ 的 \Term{borel-subgroup}{Borel 子群}）情形的推广，见 [Ja87]，II 1.20（及其中的参考文献）。}

\SourcePage{259}
\textbf{7.4. 一个例子}
作为例子，我们将确定
\[\operatorname{Hom}_{S\text{-gr.}}(\operatorname{SL}_{2,S},\operatorname{SL}_{2,S}).\]
\textbf{7.4.1。} 回顾（第 XX 讲第 5 节），$\operatorname{SL}_{2,S}$ 为由 $S$ 上满足 $ad-bc=1$ 的矩阵 $\left(\begin{smallmatrix}a&b\\c&d\end{smallmatrix}\right)$ 组成的 $S$-群\Term{scheme}{概形}。\Term{monomorphism}{单态射} $\alpha^*:\mathbb G_{m,S}\to\operatorname{SL}_{2,S}$，
\[\alpha^*(z)=\begin{pmatrix}z&0\\0&z^{-1}\end{pmatrix},\]
定义一个\Term{maximal-torus}{极大环面} $T$。$G$ 相对于 $T$ 的一个根由 $\alpha(\alpha^*(z))=z^2$ 定义，对应的\Term{monomorphism}{单态射} $p:\mathbb G_{a,S}\to\operatorname{SL}_{2,S}$ 为
\[p(x)=\begin{pmatrix}1&x\\0&1\end{pmatrix}.\]
最后，对应于 $u=p(1)$ 的 Weyl 群代表为
\[w=\begin{pmatrix}0&1\\-1&0\end{pmatrix}.\]
另回顾（第 XX 讲 6.2），$\operatorname{SL}_{2,S}$ 由 $T,p(\mathbb G_{a,S}),w$ 生成，服从关系
\[\begin{aligned}
\alpha^*(z)p(x)\alpha^*(z^{-1})&=p(z^2x),\\
w\alpha^*(z)w^{-1}&=\alpha^*(z^{-1}),\\
w^2&=\alpha^*(-1),\\(wu)^3&=1.
\end{aligned}\]
\textbf{7.4.2。} 设 $f$ 为 $\operatorname{SL}_{2,S}$ 的\Term{endomorphism}{自同态}。它首先定义\Term{homomorphism}{同态} $f\circ\alpha^*:\mathbb G_{m,S}\to\operatorname{SL}_{2,S}$。核 $\operatorname{Ker}(f\circ\alpha^*)$ 为 $\mathbb G_{m,S}$ 的\Term{closed-subgroup}{闭子群}，因此在 $S$ 上\Term{local-scheme}{局部}等于某个 $\mu_{n,S}$（$n\geq1$）或 $\mathbb G_{m,S}$。适当缩小 $S$，可假设存在 $n\in\mathbb N$ 和\Term{monomorphism}{单态射} $f':\mathbb G_{m,S}\to\operatorname{SL}_{2,S}$，使 $f\circ\alpha^*(z)=f'(z^n)$。
由\Term{monomorphism}{单态射} $\mathbb G_{m,S}\to\operatorname{SL}_{2,S}$ 的\Term{conjugacy}{共轭}性，作一个对\Term{etale}{平展}\Term{topology}{拓扑}\Term{covering}{覆盖}的底\Term{group-extension}{扩张}后，可找到 $G$ 的\Term{section}{截面} $g$，使 $f\circ\alpha^*(z)=\operatorname{int}(g)\circ\alpha^*(z^n)$。用 $\operatorname{int}(g)$ 变换 $f$，就归结为存在 $n\in\mathbb N$ 使 $f\circ\alpha^*(z)=\alpha^*(z^n)$ 的情形。

\SourcePage{260}
\textbf{7.4.3。} 现在考虑\Term{morphism}{态射} $f\circ p:\mathbb G_{a,S}\to\operatorname{SL}_{2,S}$。它满足条件
\[\alpha^*(z^n)(f\circ p)(x)\alpha^*(z^n)^{-1}=(f\circ p)(z^2x).\]
若 $n=0$，立即得出 $f\circ p$ 在 $\mathbb G_{a,S}$ 的数乘变换下不变，因而是常值的。若 $n\ne0$，可应用第 XXII 讲 4.1.9；$x\mapsto x^n$ 是 $\mathbb G_{a,S}$ 的\Term{endomorphism}{自同态}，且存在 $\lambda\in\mathbb G_{m,S}$ 使
\[f\circ p(x)=p(\lambda x^n);\]
事实上，取 $\lambda=1$ 时，此关系对 $n=0$ 也成立。再次\Term{group-extension}{扩张} $S$ 后，可找到 $z\in\mathbb G_{m,S}$ 使 $z^2=\lambda$。把 $f$ 换成 $\operatorname{int}(\alpha^*(z))\circ f$，便归结到
\[f\circ\alpha^*(z)=\alpha^*(z^n),\qquad f\circ p(x)=p(x^n).\]
\textbf{7.4.4。} 最后，必须有 $f(w)\alpha^*(z^n)f(w)^{-1}=\alpha^*(z^n)^{-1}$ 且 $(f(w)u)^3=1$。由第 XX 讲 3.8，这推出 $f(w)=w^n$。由于 $G$ 由 $T,p(\mathbb G_{a,S}),w$ 生成，便得到对任意 $\left(\begin{smallmatrix}a&b\\c&d\end{smallmatrix}\right)\in G(S')$、$S'\to S$，有
\[f\begin{pmatrix}a&b\\c&d\end{pmatrix}=\begin{pmatrix}a^n&b^n\\c^n&d^n\end{pmatrix}.\]
\textbf{7.4.5。} 综上，对于\Term{etale}{平展}\Term{topology}{拓扑}，在 $S$ 上\Term{local-scheme}{局部}，对 $\operatorname{SL}_{2,S}$ 的每个\Term{endomorphism}{自同态} $f$，可以找到 $\operatorname{SL}_{2,S}$ 的（内）\Term{automorphism}{自同构} $u$ 和整数 $n\geq0$，使 $f=u\circ F_n$，其中
\[F_n\begin{pmatrix}a&b\\c&d\end{pmatrix}=\begin{pmatrix}a^n&b^n\\c^n&d^n\end{pmatrix}.\]
注意，若 $f=u\circ F_n$，则整数 $n$ 由 $f$ 的一个\Term{fiber}{纤维} $f_s$ 唯一确定，例如由 $\operatorname{Ker}(f_s)$ 确定。因此 $n$ 是 $S$ 上的\Term{locally-constant}{局部常值}函数。
\textbf{7.4.6。} 立即推得，若 $f$ 为 $\operatorname{SL}_{2,S}$ 的\Term{endomorphism}{自同态}，则 $S$ \Term{canonical}{典范}地\Term{resolution}{分解}成既开又闭子\Term{scheme}{概形} $S_0,S_1,S_{p^n}$ 的和（其中 $p^n$ 遍历素数的所有 $>0$ 的幂），使得：

(i) $f_{S_0}$ 为零\Term{morphism}{态射}；

(ii) $f_{S_1}$ 为\Term{isomorphism}{同构}（即\Term{inner-automorphism}{内自同构}）；

(iii) $S_{p^n}$ 具有\Term{characteristic}{特征} $p$，且 $f_{S_{p^n}}$ 唯一地\Term{resolution}{分解}成 $u\circ F_p^n$，其中 $u$ 为\Term{inner-automorphism}{内自同构}，$F_p$ 为 $\operatorname{SL}_{2,\mathbb F_p}$ 的 Frobenius \Term{endomorphism}{自同态}。

\textbf{7.4.7。} 换言之，$\operatorname{Hom}_{\mathbb Z\text{-gr.}}(\operatorname{SL}_{2,\mathbb Z},\operatorname{SL}_{2,\mathbb Z})$ 是下列\Term{scheme}{概形}的和：

(i) 一个\Term{isomorphism}{同构}于 $\operatorname{Spec}(\mathbb Z)$ 的\Term{scheme}{概形}；

(ii) 一个\Term{isomorphism}{同构}于 $\operatorname{Aut}_{\mathbb Z\text{-gr.}}(\operatorname{SL}_{2,\mathbb Z})\simeq\operatorname{ad}(\operatorname{SL}_{2,\mathbb Z})$ 的\Term{scheme}{概形}；
\SourcePage{261}
(iii) 对每个素数 $p$ 和每个整数 $n>0$，一个\Term{isomorphism}{同构}于 $\operatorname{Aut}_{\mathbb F_p\text{-gr.}}(\operatorname{SL}_{2,\mathbb F_p})\simeq\operatorname{ad}(\operatorname{SL}_{2,\mathbb F_p})$ 的\Term{scheme}{概形}。

\textbf{7.4.8。} 特别地，可得：

(i) $\operatorname{Hom}_{\mathbb F_p\text{-gr.}}(\operatorname{SL}_{2,\mathbb F_p},\operatorname{SL}_{2,\mathbb F_p})$ 有无穷多个\Term{connected-component}{连通分支}，因而不\Term{quasi-compact}{拟紧}。

(ii) 若 $S$ 为混合\Term{characteristic}{特征}\Term{scheme}{概形}，则 $\operatorname{Hom}_{S\text{-gr.}}(\operatorname{SL}_{2,S},\operatorname{SL}_{2,S})$ 在 $S$ 上不\Term{flat}{平坦}。

\SGAsection{8}{附录：Hensel 环上光滑群的上同调；上同调与函子 $\prod$}
\begin{SGAresult}{命题}{8.1}
设 $S$ 为 Hensel \Term{local-scheme}{局部}\Term{scheme}{概形}，$s$ 为其\Term{closed-point}{闭点}，$G$ 为\Term{smooth}{光滑} $S$-群\Term{scheme}{概形}，并假设 $G$ 的每个\Term{finite}{有限}\Term{subset}{子集}都包含于一个\Term{affine-open}{仿射开集}中。\footnote{此最后条件实际上不需要（参见 A. Grothendieck，《Brauer 群 III》，载于《\Term{scheme}{概形}\Term{cohomology}{上同调}十讲》，North-Holland，1968，定理 11.7 和注记 11.8.3）。}则：

(i) 若 $P$ 为 $G$ 下的\Term{principal-homogeneous-bundle}{主齐性丛}，则存在\Term{finite}{有限}\Term{etale}{平展}\Term{surjectivity}{满射} $S'\to S$ 使 $P$ \Term{trivialization}{平凡化}。有 $\operatorname{Fib}(S,G)\simeq H^1_{\text{ét}}(S,G)$。

(ii) 对任意\Term{finite}{有限}\Term{etale}{平展}\Term{surjectivity}{满射} $S'\to S$，\Term{canonical}{典范}映射
\[H^1(S'/S,G)\longrightarrow H^1(S'\otimes_S\kappa(s)/\kappa(s),G_s)\]
为\Term{bijection}{双射}。

(iii) \Term{canonical}{典范}映射 $\operatorname{Fib}(S,G)\to\operatorname{Fib}(\kappa(s),G_s)$ 为\Term{bijection}{双射}。
\end{SGAresult}
\textbf{8.1.2。} 若 $K$ 为 $\kappa(s)$ 的\Term{finite-separable-extension}{有限可分扩张}，则存在\Term{finite}{有限}\Term{etale}{平展}\Term{surjectivity}{满射} $S'\to S$，使 $K\simeq S'\otimes_S\kappa(s)$。\footnote{编者注：注意 $S'$ 仍为 Hensel \Term{local-scheme}{局部}\Term{scheme}{概形}。另外，保留了原文编号：没有 8.1.1。}若 $P$ 为 $G$ 下的\Term{principal-homogeneous-bundle}{主齐性丛}，则 $P$ 在 $S$ 上\Term{smooth}{光滑}，因而 $P_s$ 在 $\kappa(s)$ 上\Term{smooth}{光滑}；故存在 $\kappa(s)$ 的\Term{finite-separable-extension}{有限可分扩张} $K$，使 $P_K$ 有\Term{section}{截面}（参见 EGA IV$_4$，17.15.10）。如上表示 $K$ 后，由“Hensel 引理”（参见 EGA IV$_4$，18.5.17）可见 $P_{S'}$ 有\Term{section}{截面}，从而证明了 (i) 的第一部分。
反之，若 $P$ 为\Term{etale}{平展}\Term{topology}{拓扑}下 $G$ 的\Term{principal-homogeneous}{主齐性}\Term{sheaf}{层}，则存在\Term{finite}{有限}\Term{etale}{平展}\Term{surjectivity}{满射} $S'\to S$ 使 $P$ \Term{trivialization}{平凡化}（事实上，一个 Hensel \Term{local-scheme}{局部}\Term{scheme}{概形}的任意\Term{etale}{平展}\Term{covering-family}{覆盖族}，都被一个仅含有所需\Term{form}{形式}\Term{morphism}{态射} $S'\to S$ 的族\Term{refinement}{加细}）。
由对 $G$ 所作的\Term{descent}{下降}假设，$P$ \Term{representable}{可表}（SGA 1，第 VIII 讲 7.6）；这完成 (i) 的证明，并表明 (ii) 推出 (iii)。所以只须证明 (ii)。

\SourcePage{262}
\textbf{8.1.3。} (ii) 中的映射是\Term{injectivity}{单射}：设 $P,Q$ 为两个由 $S'$ \Term{trivialization}{平凡化}的 $G$ 下的\Term{principal-homogeneous-bundle}{主齐性丛}。考虑 $S$-群\Term{sheaf}{层} $H=\operatorname{Isom}_{G\text{-丛}}(P,Q)$；由于 $H_{S'}$ \Term{isomorphism}{同构}于 $G_{S'}$，由上面对 $G$ 的第二个假设，$H$ \Term{representable}{可表}。若 $H(\kappa(s))\ne\varnothing$，则由 Hensel 引理有 $H(S)\ne\varnothing$，所以 $P,Q$ \Term{isomorphism}{同构}。

\textbf{8.1.4。} 最后证明 (ii) 中的映射\Term{surjectivity}{满射}，亦即\Term{canonical}{典范}映射
\[Z^1(S'/S,G)\longrightarrow Z^1(S'\otimes_S\kappa(s)/\kappa(s),G_s)\]
\Term{surjectivity}{满射}。设 $Z^1$ 为如下定义的 $S$-\Term{functor}{函子}：
\[Z^1(T)=Z^1(S'\times_ST/T,G_T);\]
上述映射等同于映射 $Z^1(S)\to Z^1(\kappa(s))$；再次应用 Hensel 引理，只须证明 $Z^1$ 可由\Term{smooth}{光滑} $S$-\Term{scheme}{概形}表示。

\textbf{8.1.5。} 证明 $Z^1$ 可由\Term{locally-finite-presentation}{局部有限表示} $S$-\Term{scheme}{概形}表示。设 $C^i$（$i=0,1,\ldots$）为如下定义的 $S$-\Term{functor}{函子}：
\[C^i(T)=C^i(S'\times_ST/T,G_T),\]
即
\[C^i(T)=G((S'\times_ST/T)^{i+1})=G((S'/S)^{i+1}\times_ST),\]
或写成 $C^i=\operatorname{Hom}_S((S'/S)^{i+1},G)$。由于 $Z^1$ 由 $C^1,C^2$ 通过\Term{fiber-product}{纤维积}得到，只须证明 $C^i$（$i=1,2$）可由\Term{locally-finite-presentation}{局部有限表示} $S$-\Term{scheme}{概形}表示。

\textbf{8.1.6。} 若 $T\to S$ 为使 $S'$ \Term{resolution}{分解}的\Term{finite}{有限}\Term{etale}{平展}\Term{surjectivity}{满射}，则
\[C^i\times_ST=\operatorname{Hom}_T((S'\times_ST/T)^{i+1},G_T)\]
可由 $n$ 个 $G_T$ 的副本的乘积表示，其中 $n$ 为 $(S'/S)^{i+1}$ 的次数。再次应用对 $G$ 的假设，就得到 $C^i$ 确实可由\Term{locally-finite-presentation}{局部有限表示} $S$-\Term{scheme}{概形}表示（SGA 1，第 VIII 讲，上引处）。

\textbf{8.1.7。} 为证明 $Z^1$ \Term{smooth}{光滑}，现在根据定义，须证明：若 $T$ 为\Term{affine}{仿射} $S$-\Term{scheme}{概形}，$T_0$ 为一个\Term{square-zero-ideal}{平方零理想} $\mathscr J$ 所定义的闭子\Term{scheme}{概形}，则\Term{canonical}{典范}映射 $Z^1(T)\to Z^1(T_0)$ \Term{surjectivity}{满射}。由于 $G$ \Term{smooth}{光滑}，\Term{canonical}{典范}映射 $G(T)\to G(T_0)$ \Term{surjectivity}{满射}，
\SourcePage{263}
所以只须证明\Term{canonical}{典范}映射
\[H^1(S'\times_ST/T,G_T)\longrightarrow H^1(S'\times_ST_0/T_0,G_{T_0})\]
\Term{bijection}{双射}。稍改记号并推广假设后，现在只须证明：
\begin{SGAresult}{引理}{8.1.8}
设 $S,S'$ 为两个\Term{affine}{仿射}\Term{scheme}{概形}，$S'\to S$ 为\Term{faithfully-flat}{忠实平坦}\Term{morphism}{态射}，$\mathscr J$ 为 $S$ 上的\Term{square-zero-ideal}{平方零理想}，$S_0$ 为其定义的闭子\Term{scheme}{概形}，$G$ 为\Term{smooth}{光滑} $S$-群。置 $S'_0=S'\times_SS_0$、$G_0=G\times_SS_0$。\Term{canonical}{典范}映射
\[H^1(S'/S,G)\longrightarrow H^1(S'_0/S_0,G_0)\]
为\Term{bijection}{双射}。
\end{SGAresult}
\textbf{注记 8.1.9。} 若假设 $G$ 交换，则同样断言对所有 $H^i$（$i>0$）成立，证明类似。

\begin{proof}
设 $\mathscr M_0$ 为\Term{quasi-coherent}{拟凝聚} $\mathscr O_{S_0}$-模 $\mathscr M_0=\operatorname{Lie}(G_0/S_0)\otimes_{\mathscr O_{S_0}}\mathscr J$。对每个 $S_0$-\Term{prescheme}{预概形} $T_0$，置 $\mathbf M_0(T_0)=H^0(T_0,\mathscr M_0\otimes_{\mathscr O_{S_0}}\mathscr O_{T_0})$，并设
\[\mathbf M=\prod_{S_0/S}\mathbf M_0,\qquad\overline G=\prod_{S_0/S}G_0\]
为由 $\mathbf M(T)=\mathbf M_0(T_0)$ 和 $\overline G(T)=G_0(T_0)$ 定义的 $S$-\Term{group-functor}{群函子}，其中 $T_0=T\times_SS_0$。根据第 III 讲 0.9 和 (0.6.2)，对每个\Term{affine}{仿射} $S$-\Term{scheme}{概形} $T$，有关于 $T$ \Term{functorial}{函子性}的\Term{exact-sequence}{正合列}
\[1\longrightarrow\mathbf M(T)\longrightarrow G(T)\longrightarrow\overline G(T)\longrightarrow1.\]
须研究的映射是
\[H^1(S'/S,G)\longrightarrow H^1(S'_0/S_0,G_0)=H^1(S'/S,\overline G).\]
先假设 $G$ 交换。此时有\Term{cohomology}{上同调}\Term{exact-sequence}{正合列}
\[H^1(S'/S,\mathbf M)\longrightarrow H^1(S'/S,G)\longrightarrow H^1(S'/S,\overline G)\longrightarrow H^2(S'/S,\mathbf M);\]
而
\[H^i(S'/S,\mathbf M)=H^i(S'_0/S_0,\mathbf M_0)=H^i(S'_0/S_0,\mathscr M_0),\]
且已知（TDTE I，B，引理 1.1），$i\ne0$ 时 $H^i(S'_0/S_0,\mathscr M_0)=0$。
若 $G$ 不交换，就必须使用\Term{nonabelian-cohomology}{非阿贝尔上同调}\Term{exact-sequence}{正合列}。若 $u\in Z^1(S'/S,G)$，已知 $H^1(S'/S,G)$ 中与 $u$ 的类在 $H^1(S'/S,\overline G)$ 中有相同像的元素，都位于相应的上边缘映射
\[H^1(S'/S,\mathbf M_u)\longrightarrow H^1(S'/S,G)\]
的像中，其中 $\mathbf M_u$ 为“被 $u$ \Term{twisting}{扭曲}”的 $S$-\Term{functor}{函子} $\mathbf M$。同样，若 $v\in Z^1(S'/S,\overline G)$，则存在“上边缘”
\[\Delta(v)\in H^2(S'/S,\mathbf M_v),\]
\SourcePage{264}
其中 $\mathbf M_v$ 为“被 $v$ \Term{twisting}{扭曲}”的 $S$-\Term{functor}{函子} $\mathbf M$，且 $\Delta(v)=0$ 当且仅当 $v$ 的类在 $H^1(S'/S,G)$ 的像中。只须证明
\[H^1(S'/S,\mathbf M_u)=H^2(S'/S,\mathbf M_v)=0.\]
回顾（第 III 讲 0.8），上述\Term{exact-sequence}{正合列}定义的 $G$ 在 $\mathbf M$ 上的\Term{action}{作用}，正是\Term{functorial}{函子性}地由 $G_0$ 的\Term{adjoint-representation}{伴随表示}
\[\operatorname{ad}:G_0\longrightarrow\operatorname{Aut}_{\mathscr O_{S_0}}(\operatorname{Lie}(G_0/S_0))\]
得到的\Term{action}{作用}。因而元素 $u$（相应地，$v$）通过 $\operatorname{Lie}(G_0/S_0)$ 的一个 $S'_0\times_{S_0}S'_0$-\Term{automorphism}{自同构}\Term{action}{作用}于 $\mathbf M_{S'\times_SS'}$。由于 $u$（相应地，$v$）为余圈，此\Term{automorphism}{自同构}是\Term{descent}{下降}数据；记所得的\Term{quasi-coherent}{拟凝聚} $\mathscr O_{S_0}$-模为 $\mathscr L_u$（相应地，$\mathscr L_v$）。立即可验证，对 $T\to S$，有
\[\mathbf M_u(T)=H^0(T_0,\mathscr L_u\otimes_{\mathscr O_{S_0}}\mathscr J\otimes_{\mathscr O_{S_0}}\mathscr O_{T_0}),\]
把 $u$ 换成 $v$ 也有相同关系。因此
\[\begin{aligned}
H^1(S'/S,\mathbf M_u)&=H^1(S'_0/S_0,\mathscr L_u\otimes\mathscr J),\\
H^2(S'/S,\mathbf M_v)&=H^2(S'_0/S_0,\mathscr L_v\otimes\mathscr J),
\end{aligned}\]
由已用过的结果，两者确实均为零。
\end{proof}

\begin{SGAresult}{命题}{8.2}
设 $\mathscr C$ 为有\Term{fiber-product}{纤维积}的\Term{category}{范畴}，配备一个比\Term{canonical-topology}{典范拓扑}更粗的\Term{topology}{拓扑}，$S'\to S$ 为 $\mathscr C$ 的\Term{morphism}{态射}，$G'$ 为 $S'$-群\Term{sheaf}{层}，$G$ 为 $S$-群\Term{sheaf}{层} $\prod_{S'/S}G'$。设 $H^1_S(S',G')\subset H^1(S',G')$ 为如下\Term{principal-homogeneous}{主齐性}\Term{sheaf}{层}的类所成的\Term{set}{集合}：它们在 $G'$ 下，并被某个由 $S$ 的适当\Term{covering}{覆盖}\Term{sieve}{筛}\Term{base-change}{基变换}而得的 $S'$ 的\Term{sieve}{筛}\Term{trivialization}{平凡化}。由\Term{functor}{函子}
\[P\longmapsto P\times_SS'\]
定义的\Term{canonical}{典范}映射 $H^1(S,G)\to H^1(S',G')$ 诱导\Term{bijection}{双射}
\[H^1(S,G)\xrightarrow{\sim}H^1_S(S',G');\]
其逆\Term{bijection}{双射}由\Term{functor}{函子} $P'\mapsto\prod_{S'/S}P'$ 定义。
\end{SGAresult}
对 $\mathscr C_{/S}$ 的每个\Term{object}{对象} $X$，按定义有关于 $X$ \Term{functorial}{函子性}的\Term{isomorphism}{同构}
\[\operatorname{Hom}_S(X,G)\xrightarrow{\sim}\operatorname{Hom}_{S'}(X\times_SS',G').\]
所以对每个 $S$-\Term{object}{对象} $T$，有关于 $T$ \Term{functorial}{函子性}的\Term{bijection}{双射}
\[H^1(T/S,G)\xrightarrow{\sim}H^1(T'/S',G').\]
再把单个\Term{morphism}{态射} $T\to S$ 换成 $S$ 的任意\Term{covering-family}{覆盖族}，并取\Term{inductive-limit}{归纳极限}，就得到陈述的第一部分。第二部分由此容易推出。
\begin{SGAresult}{引理}{8.3}
在 8.2 的条件下，断言 $H^1_S(S',G')=H^1(S',G')$ 在 $S$ 上是\Term{local-scheme}{局部}的：设存在\Term{covering-family}{覆盖族} $\{S_i\to S\}$，使对每个 $i$ 都有
\[H^1_{S_i}(S'\times_SS_i,G')=H^1(S'\times_SS_i,G').\]
则 $H^1_S(S',G')=H^1(S',G')$。
\end{SGAresult}
\SourcePage{265}
事实上，设 $P'$ 为 $G'$ 下的\Term{principal-homogeneous}{主齐性}\Term{sheaf}{层}。置
\[P'_i=P\mathop{\times}_{S'}(S'\times_SS_i);\]
由假设，存在\Term{covering-family}{覆盖族} $\{S_{ij}\to S_i\}$，使对每个 $j$，$P'\times_{S'}(S'\times_SS_{ij})$ 有\Term{section}{截面}。但 $\{S_{ij}\to S\}$ 是 $S$ 的\Term{covering-family}{覆盖族}，且 $P'$ 确实被此族经\Term{base-change}{基变换}所得的 $S'$ 的\Term{covering-family}{覆盖族}\Term{trivialization}{平凡化}。
\begin{SGAresult}{命题}{8.4}
设 $S'\to S$ 为\Term{scheme}{概形}的\Term{finite}{有限}\Term{etale}{平展}\Term{morphism}{态射}。设 $G'$ 为 $S'$-群\Term{sheaf}{层}，$G$ 为 $S$-群\Term{sheaf}{层} $\prod_{S'/S}G'$。对于\Term{etale}{平展}\Term{topology}{拓扑}（相应地，\Term{locally-finite}{局部有限}\Term{etale}{平展}\Term{topology}{拓扑}，或 (fpqc)），\Term{functor}{函子}
\[\begin{aligned}
P&\longmapsto P\times_SS',\\
\prod_{S'/S}P'&\longleftarrow P'
\end{aligned}\]
诱导互逆的\Term{bijection}{双射}
\[H^1(S,G)\simeq H^1(S',G').\]
\end{SGAresult}
由 8.2，只须证明 $H^1_S(S',G')=H^1(S',G')$。由 8.3，只须对\Term{locally-finite}{局部有限}\Term{etale}{平展}\Term{topology}{拓扑}\Term{local-scheme}{局部}地证明；因而可假设 $S'$ 为\Term{finite}{有限}多个 $S$ 的副本的\Term{direct-sum}{直和}，即 $I_S$，其中 $I$ 为适当的\Term{finite}{有限}\Term{set}{集合}。于是 $G'$ 由 $S$ 上的\Term{sheaf}{层}族 $(G_i)_{i\in I}$ 给出，且
\[H^1(S',G')\simeq\prod_{i\in I}H^1(S,G_i).\]
另一方面，$H^1(S,G)\simeq\prod_{i\in I}H^1(S,G_i)$，故由 8.2，$H^1_S(S',G')=H^1(S',G')$。证毕。
\textbf{注记 8.5。} 可以通过下列\Term{exact-sequence}{正合列}解释 8.2 和 8.3（$f$ 为给定\Term{morphism}{态射} $S'\to S$）：
\[1\longrightarrow H^1(S,f_*(G'))\longrightarrow H^1(S',G')\longrightarrow H^0(S,R^1f_*(G')).\]
在交换情形，此\Term{exact-sequence}{正合列}由 \Term{leray-spectral-sequence}{Leray 谱序列}得到；它在非交换情形仍成立（参见 Giraud 的学位论文\footnote{编者注：见 [Gi71]，V 3.1.4。}）。
在此\Term{form}{形式}下，可见若只假设 $f$ \Term{finite}{有限}，乃至仅为\Term{integral-morphism}{整态射}，只要所取\Term{topology}{拓扑}为\Term{etale}{平展}\Term{topology}{拓扑}，结果仍成立；因为由 SGA 4，第 VIII 讲 5.3，\footnote{编者注：这一引用也指向 [Gi71]，V 3.1.4。}对任意 $G'$ 都有
\[R^1f_*(G')=\text{终层}.\]
另一方面，若取 (fpqc) 或 (fppf) 这样的\Term{topology}{拓扑}，结果便不成立，即使 $S=\operatorname{Spec}(k)$，$k$ 为\Term{characteristic}{特征} $p\ne0$ 的\Term{algebraically-closed-field}{代数闭域}，$S'=\operatorname{Spec}(k[t]/t^2)$，$G'=\mu_p$ 或 $\alpha_p$，也是如此。

\SourcePage{266}
同样，若在 8.2 中去掉 $f$ \Term{finite}{有限}的假设，则即使对于\Term{etale}{平展}\Term{topology}{拓扑}，结论也不成立；取 $f$ 为一个\Term{open-immersion}{开浸入}即可看出。例如，$S=\operatorname{Spec}(V)$，$V$ 为\Term{residue-field}{剩余域}\Term{algebraically-closed}{代数闭}的\Term{complete-discrete-valuation-ring}{完备离散赋值环}，$S'$ 为\Term{generic-point}{一般点}处诱导的\Term{open-subset}{开集}，$G'$ 为\Term{constant-group}{常值群} $(\mathbb Z/n\mathbb Z)_{S'}$，其中 $n$ 与 $V$ 的\Term{residue-characteristic}{剩余特征}互素，则 $H^1(S,G)=0$，$H^1(S',G')\ne0$。再把 $S'$ 换成 $S\amalg S'$，即可得到类似的例子，其中 $S'\to S$ 为\Term{etale}{平展}\Term{surjectivity}{满射}，因而是所考虑\Term{topology}{拓扑}的\Term{covering}{覆盖}。

\textbf{参考文献}
\begin{description}
\item[{[BLie]}] N. Bourbaki，《Lie 群与 \Term{lie-algebra}{Lie 代数}》，第 I 章，Hermann，1960。
\item[{[Ch51]}] C. Chevalley，《Lie 群理论》，第 II 卷《\Term{algebraic-group}{代数群}》，Hermann，1951。
\item[{[Di57]}] J. Dieudonné，《\Term{characteristic}{特征} $p>0$ 的域上的 Lie 群与 Lie 超\Term{algebra}{代数}，VI》，\textit{Amer. J. Math.} 79 (1957)，第 2 期，331--388。\footnote{编者注：在原文列出的这三项参考文献之后，添加了以下各项。}
\item[{[Bo91]}] A. Borel，《线性\Term{algebraic-group}{代数群}》，第 2 版，Springer-Verlag，1991。
\item[{[Ch05]}] C. Chevalley，《\Term{semisimple}{半单}\Term{algebraic-group}{代数群}的分类》（P. Cartier、A. Grothendieck、M. Lazard 合作），\textit{Collected Works}，第 3 卷，Springer，2005。
\item[{[DG70]}] M. Demazure、P. Gabriel，《\Term{algebraic-group}{代数群}》，Masson \& North-Holland，1970。
\item[{[Gi71]}] J. Giraud，《\Term{nonabelian-cohomology}{非阿贝尔上同调}》，Springer-Verlag，1971。
\item[{[Ja87]}] J. C. Jantzen，《\Term{algebraic-group}{代数群}的表示》，Academic Press，1987；第 2 版，Amer. Math. Soc.，2003。
\item[{[Jou83]}] J.-P. Jouanolou，《Bertini 定理及其应用》，Birkhäuser，1983。
\item[{[Ta75]}] M. Takeuchi，《\Term{affine-algebraic-group}{仿射代数群}的\Term{covering}{覆盖}与超\Term{algebra}{代数}》，\textit{Trans. Amer. Math. Soc.} 211 (1975)，249--275。
\item[{[Ta76]}] M. Takeuchi，《\Term{affine-algebraic-group}{仿射代数群}的\Term{covering}{覆盖}与超\Term{algebra}{代数} II》，\textit{J. Fac. Sci. Univ. Tokyo Sect. IA Math.} 23 (1976)，419--434。
\end{description}
